\documentclass[11pt,a4paper]{article}
\usepackage[a4paper,margin=0.95in]{geometry}
\usepackage[T1]{fontenc}
\usepackage{lmodern,microtype}
\usepackage{amsmath,amssymb,amsthm,mathtools}
\usepackage{booktabs,array,tabularx,longtable,multirow,float}
\usepackage{enumitem,xcolor,url,listings,needspace}
\usepackage[hidelinks]{hyperref}
\usepackage{xr-hyper}
\usepackage[nameinlink,noabbrev,capitalise]{cleveref}
\usepackage{fancyhdr}

\externaldocument[][nocite]{output/pdf/QRUOV_supplement}[QRUOV_supplement.pdf]

\definecolor{deepblue}{RGB}{28,69,125}
\hypersetup{colorlinks=true,linkcolor=deepblue,citecolor=deepblue,urlcolor=deepblue,
 pdftitle={New Attacks on QR-UOV Across Vinegar Dimensions: Direct Forgery and Equivalent-Key Recovery},
 pdfauthor={Justin Thaler, Kai Zhang, Scott Kominers, Quang Dao},
 pdfsubject={Projective direct forgery, verified equivalent-key recovery, and parameter implications for QR-UOV},
 pdfkeywords={QR-UOV,UOV,projective direct forgery,equivalent-key recovery,Chevalley-Warning,full-rank roots,Karatsuba,symbolic equation solving,exact verification}}

\theoremstyle{definition}
\newtheorem{theorem}{Theorem}[section]
\newtheorem{lemma}[theorem]{Lemma}
\newtheorem{proposition}[theorem]{Proposition}
\newtheorem{corollary}[theorem]{Corollary}
\newtheorem{definition}[theorem]{Definition}
\theoremstyle{remark}
\newtheorem{remark}[theorem]{Remark}

\newcommand{\F}{\mathbb F}
\newcommand{\PP}{\mathbb P}
\newcommand{\cO}{\mathcal O}
\newcommand{\cB}{\mathcal B}
\newcommand{\cX}{\mathcal X}
\newcommand{\cR}{\mathcal R}
\newcommand{\cU}{\mathcal U}
\newcommand{\cV}{\mathcal V}
\newcommand{\cW}{\mathcal W}
\newcommand{\cG}{\mathcal G}
\newcommand{\cE}{\mathcal E}
\newcommand{\Span}{\operatorname{span}}
\newcommand{\rank}{\operatorname{rank}}
\newcommand{\codim}{\operatorname{codim}}
\newcommand{\cdeg}{\operatorname{cdeg}}
\newcommand{\Sym}{\operatorname{Sym}}
\newcommand{\disc}{\mathrm{disc}}
\newcommand{\bad}{\mathrm{bad}}
\newcommand{\oil}{\mathrm{oil}}
\newcommand{\off}{\mathrm{off}}
\newcommand{\spanbad}{\mathrm{span}}
\newcommand{\QRUOV}{\textsf{QR-UOV}}
\newcommand{\Solve}{\operatorname{Solve}}
\newcommand{\Adv}{\operatorname{Adv}}
\newcommand{\trans}{\mathsf T}
\newcommand{\softO}{\widetilde O}
\newcommand{\Mmul}{\mathsf M}
\newcommand{\Tr}{\operatorname{Tr}}
\newcommand{\Gmul}{\mathsf G_{\rm mul}}
\newcommand{\proofstep}[1]{\par\smallskip\noindent\textbf{#1}\enspace}
\newcommand{\ord}{\operatorname{ord}}
\newcommand{\M}{\mathsf M}

\setlist[itemize]{leftmargin=*,itemsep=0.10em,topsep=0.20em}
\setlist[enumerate]{leftmargin=*,itemsep=0.12em,topsep=0.20em}
\setlength{\parindent}{1.2em}
\setlength{\parskip}{0.20em}
\linespread{1.04}
\setlength{\emergencystretch}{2em}
\raggedbottom
\pagestyle{fancy}
\fancyhf{}
\fancyfoot[C]{\thepage}
\renewcommand{\headrulewidth}{0pt}

\title{\bfseries New Attacks on \QRUOV{} Across Vinegar Dimensions:\ Direct Forgery and Equivalent-Key Recovery}
\author{Justin Thaler\thanks{Georgetown University and a16z crypto research.}
\and Kai Zhang\thanks{MIT.}
\and Scott Kominers\thanks{Harvard University and a16z crypto research.}
\and Quang Dao\thanks{Carnegie Mellon University and LayerZero Labs.}}
\date{31 July 2026}

\begin{document}
\maketitle

\begin{abstract}
\QRUOV{} is a structured variant of Unbalanced Oil and Vinegar (UOV) in the
third round of NIST's Additional Digital Signatures process.  It keeps UOV's
short signatures and simple signing procedure, but compresses its large public
key by representing blocks of the public quadratic equations over a
degree-three extension field.  Previous analyses study forgery and key
recovery separately.  This leaves open whether changing the number of vinegar
variables can make both attacks harder.

We prove that it cannot.  Let $m$ be the number of public equations and $\ell$
the extension degree.  For every number $V$ of extension-field vinegar
variables and every nonzero verification target, one attack reduces to solving
at most $m-\ell-1$ nonlinear equations in the same number of variables.  When
$V$ is small, this is an equivalent-key recovery system.  When $V$ is large,
it is a direct-forgery system.  Changing $V$ therefore only shifts the system
between key recovery and direct forgery.  In \QRUOV{},
$\ell=3$, so the bound is $m-4$, or $50$, $74$, and $101$ variables for the
three recommended parameter sets.

We estimate the work by adapting a symbolic equation solver to \QRUOV{}'s base
field and applying an existing Boolean-gate cost model.  It gives work factors
of $2^{140.384}$, $2^{191.027}$, and $2^{247.125}$ gates, below the claimed
security levels of $143$, $207$, and $272$ bits.  On smaller examples, our
direct-forgery implementation runs from public input to a forged signature
accepted by the verifier.  This confirms that the attack steps work together,
not that the full-parameter attacks are practical.  The estimates depend on
extrapolating the solver's behavior, the assumed cost of its polynomial
arithmetic, and a pseudorandomness model for the expanded public key.  They
also require an impractical amount of memory.
\end{abstract}

\section{Introduction}\label{sec:intro}

NIST opened its Additional Digital Signatures process to broaden the set of
post-quantum signatures beyond the structured-lattice schemes selected in its
first standardization process.  After the second round, NIST advanced nine
candidates, including \QRUOV{}, to a third round of public analysis
~\cite{NISTRound3,NISTIR8610}.  \QRUOV{} retains the short signatures and
simple signing procedure of Unbalanced Oil and Vinegar (UOV) while reducing
UOV's main efficiency cost: its large public key
~\cite{QRUOVOriginal,QRUOVSpec}.

UOV divides the secret coordinates into vinegar variables and oil variables.
Its central quadratic equations contain no product of two oil variables.
After choosing the vinegar variables, the signer can therefore solve a linear
system for the oil variables.  A secret change of coordinates hides this
division in the public equations.  Using more vinegar variables can make the
hidden oil space harder to recover, but it also gives a direct forger more
variables with which to solve the public equations
~\cite{KipnisShamir1998,KPG1999,BraekenWolfPreneel2005}.

We next describe \QRUOV{} in the notation used throughout the paper.  Fix an
odd prime power $q$ and an irreducible polynomial $f\in\F_q[X]$ of degree
$\ell$.  Write
\begin{equation}\label{eq:parameter-stack}
 K:=\F_q[X]/(f)\cong\F_{q^\ell},\qquad
 v=\ell V,\qquad m=\ell O,\qquad
 N=V+O,\qquad n=v+m=\ell N.
\end{equation}
The verifier evaluates $m$ quadratic forms in $n$ variables over $\F_q$.
The quotient representation groups these variables into $V$ vinegar blocks
and $O$ oil blocks over $K$.  The number of output forms remains $m$, rather
than falling to $O=m/\ell$.

For each output $i\in[m]$, the secret central matrix and the secret input
transformation have the following form over $K$:
\begin{equation}\label{eq:central}
 F_i=\begin{pmatrix}A_i&B_i\\B_i^{\trans}&0\end{pmatrix},
 \qquad
 S_G=\begin{pmatrix}I_V&G\\0&I_O\end{pmatrix},
 \qquad
 P_i=S_G^{\trans}F_iS_G
     =\begin{pmatrix}A_i&E_i\\E_i^{\trans}&C_i\end{pmatrix}.
\end{equation}
Here $G\in K^{V\times O}$ is secret,
$E_i=A_iG+B_i$, and block multiplication gives
\begin{equation}\label{eq:graphid}
 C_i=G^{\trans}E_i+E_i^{\trans}G-G^{\trans}A_iG.
\end{equation}
The zero lower-right block of $F_i$ is the UOV condition.  For extension-field
vinegar and oil coordinates $(z,c)\in K^V\times K^O$, it gives
\begin{equation}\label{eq:central-evaluation}
 F_i(z,c)=z^{\trans}A_iz+2z^{\trans}B_ic.
\end{equation}
Once $z$ is fixed, every central quadratic form is linear in $c$.

The same matrices identify the hidden oil space in the public coordinates.
It is the graph
\begin{equation}\label{eq:graph}
 \cO=S_G^{-1}(\{0\}^V\times K^O)
     =\left\{\binom{-Gc}{c}:c\in K^O\right\}.
\end{equation}
Every public form vanishes on every pair of vectors from this space.  In
matrix notation, this means
\begin{equation}\label{eq:isotropy}
 u^{\trans}P_iw=0
 \qquad(u,w\in\cO,\ 1\le i\le m).
\end{equation}
Conversely, a matrix $G'$ defines a common oil graph for the public tuple
exactly when it satisfies \eqref{eq:graphid} with $G'$ in place of $G$.
Thus \eqref{eq:graphid} is the condition that any recovered oil graph must
pass.

We now convert the extension-field matrices into the base-field matrices used
by the signer and verifier.  The regular representation
$\Phi_f:K\to\F_q^{\ell\times\ell}$ maps an element $g\in K$ to the matrix for
multiplication by $g$ on its $\ell$ base-field coefficients.  The map
$\Phi_f^{(a,b)}$ applies this conversion to every block of a matrix in
$K^{a\times b}$.  The specification also supplies a matrix $W_f$ that makes
$W_f\Phi_f(g)$ symmetric for every $g\in K$.  Let $W_f^{(a)}$ contain $a$
copies of $W_f$ on its diagonal.  The specification's conversion, which we
call the \emph{pull-back}, gives the following base-field matrices:
\begin{equation}\label{eq:official-pullback}
 \widehat P_i=W_f^{(N)}\Phi_f^{(N,N)}(P_i),\qquad
 \widehat F_i=W_f^{(N)}\Phi_f^{(N,N)}(F_i),\qquad
 \widehat S_G=\Phi_f^{(N,N)}(S_G).
\end{equation}
The pull-back preserves the secret factorization, so
\begin{equation}\label{eq:base-factorization}
 \begin{aligned}
 \widehat P_i&=\widehat S_G^{\trans}\widehat F_i\widehat S_G,\\
 \widehat P(x)&:=(x^{\trans}\widehat P_1x,\ldots,
 x^{\trans}\widehat P_mx),\\
 \widehat F(z)&:=(z^{\trans}\widehat F_1z,\ldots,
 z^{\trans}\widehat F_mz),
 \qquad \widehat P(x)=\widehat F(\widehat S_Gx).
 \end{aligned}
\end{equation}
A general $\ell\times\ell$ base-field block has $\ell^2$ entries, while
$\Phi_f(g)$ is determined by the $\ell$ coefficients of $g\in K$.  The
official public seed expands the blocks $A_i$ and $E_i$, the public key stores
the remaining blocks $C_i$, and the secret seed expands $G$.  This is how
\QRUOV{} compresses the dense UOV verification map without changing the map
that the verifier evaluates~\cite{QRUOVOriginal,QRUOVSpec}.

The signing equation will also be the interface for both attacks.  For a
message $M$ and salt $r$, define the verification target
\begin{equation}\label{eq:verification-target}
 \mu=\operatorname{SHAKE256}(\mathsf{seed}_{\rm pk}
       \mathbin\|\operatorname{BytesToBits}(M),512),
 \qquad
 y(M,r)=\operatorname{Hash}(\mu,r)\in\F_q^m.
\end{equation}
For a vinegar assignment $a\in\F_q^v$, define the oil coefficient matrix by
\begin{equation}\label{eq:signing-matrix}
 \mathsf L_{\widehat F}(a)\in\F_q^{m\times m},
 \qquad
 \bigl(\mathsf L_{\widehat F}(a)o\bigr)_i
 :=\widehat F_i(a,o)-\widehat F_i(a,0).
\end{equation}
The right-hand side is linear in $o\in\F_q^m$ because the oil block of
$\widehat F_i$ is zero.  The signer chooses $a$ and then samples $r$ until
\begin{equation}\label{eq:honest-oil-system}
 \mathsf L_{\widehat F}(a)o
 =y(M,r)-\widehat F(a,0)
\end{equation}
has a solution~\cite[Sec.~4.2]{QRUOVSpec}.  It returns the signature $(r,x)$
with $x=\widehat S_G^{-1}(a,o)$.  The verifier's acceptance condition is
~\cite[Sec.~4.3]{QRUOVSpec}
\begin{equation}\label{eq:verification-equation}
 \widehat P(x)=y(M,r).
\end{equation}

The verifier equation exposes the two attack goals studied here.  Direct
forgery fixes
a nonzero target $y$ and finds one $x$ satisfying $\widehat P(x)=y$.
Equivalent-key recovery finds any $G'$ satisfying \eqref{eq:graphid}, pulls
the resulting central forms back to $\F_q$, and finds an $a$ for which
$\mathsf L_{\widehat F}(a)$ is invertible.  The same linear system then signs
every verification target.  Increasing the vinegar count can frustrate the
second goal, but it makes the public system more underdetermined and can help
the first.

The public key still represents $m$ dense quadratic forms, each with
quadratically many coefficients.  Field lifting and block-anti-circulant UOV
are earlier approaches to compressing these coefficients
~\cite{FieldLifting,BACUOV}; the latter was broken by decomposing the reducible
quotient ring behind its public blocks~\cite{BACUOVAttack}.  \QRUOV{} avoids
that decomposition by choosing $f$ irreducible, but the extension-field tuple
$(P_1,\ldots,P_m)$ remains public after expansion.  This tuple underlies the
QR-UOV pull-back and lifting attacks, as well as attacks that search a public
linear space for low-rank matrices
~\cite{QRUOVOriginal,FurueIkematsu}.

Existing work studies the two consequences of changing $V$ separately.
Direct attacks reduce the underdetermined public equations before invoking a
polynomial solver
~\cite{ThomaeWolf2012,FurueUnderdetermined2021,Hashimoto2023,JustGuess}.
Structural attacks instead search for the hidden oil space through
reconciliation, intersection, or MinRank equations
~\cite{DingEtAl2008,Beullens,FurueIkematsu,Suzuki}.  P\'ebereau showed that one
suitable oil vector can determine an equivalent UOV key and made the parameter
tension explicit: more variables can impede key recovery while helping
forgery~\cite{Pebereau}.  The original \QRUOV{} analysis similarly chooses the
vinegar count mainly against structural attacks and the output count against
direct forgery~\cite{QRUOVOriginal}.  The specification and the team's latest
security response evaluate these attacks at fixed recommended parameter sets
~\cite{QRUOVSpec,QRUOVResponse2026}.  They do not give one bound covering every
integer $V$ while $m$ and $\ell$ remain fixed.

We call a nonlinear system \emph{square} when it has the same number of
equations and variables.  We compare the two attack families by counting the
variables in the square system passed to the root solver.  This brings us to
the question that governs the paper:
\begin{quote}
\noindent\textbf{Question.}
\emph{Can \QRUOV{} force both attack reductions to use square nonlinear
systems with more than $m-4$ variables by changing only the vinegar count
while keeping the output dimension and quotient degree fixed?}
\end{quote}

\subsection{Our results and technical overview}\label{sec:intro-results}

We answer this question negatively.  Our main result is the following bound
for a fixed number $m$ of public equations:
\begin{quote}
\centering
\textbf{\emph{For every vinegar count and nonzero verification target,
either direct forgery or equivalent-key recovery reduces to a square nonlinear
system in \mbox{at most~$m-4$ variables}.}}
\end{quote}
The statement, formalized in \cref{thm:security-scissors}, concerns the
shared nonlinear system.  It does not assert that total running time or
memory depends only on this dimension.

\paragraph{Direct forgery.}
For a fixed nonzero target, we first change the public output coordinates so
that all but one target coordinate are zero.  We then use
Chevalley--Warning to construct an $(m-3)$-dimensional subspace on which three
of the corresponding quadratic forms vanish.  The remaining zero-target
equations are homogeneous, so multiplying a solution by a nonzero scalar does
not change them.  Fixing this common scale leaves the same number of equations
and variables, namely $m-4$.  The dimensions are $50$, $74$, and $101$ for
the three recommended parameter sets
(Theorem~\ref{thm:verified-direct-forgery} and
Remark~\ref{rem:direct-prior-art}).  The attack uses only public randomness,
and the verifier checks every returned vector against the original public map
(\cref{sec:direct}).

\paragraph{Equivalent-key recovery.}
In extension-field coordinates, the hidden oil space is the graph of a linear
map $G$.  Fixing a public oil vector $c$ produces $m$ public quadratic
equations in $V$ variables.  The secret graph guarantees that $Gc$ is a root,
although the attacker does not know this root in advance.  If the matrix of
first derivatives has full rank at a recovered root, public linear equations
determine all of $G$.  The pull-back and signing checks then certify another
secret key that signs every target accepted by the verifier
(\cref{thm:verified-directional-recovery}).  P\'ebereau established the
equivalent-key endpoint from one suitable oil vector~\cite{Pebereau}; our
contribution is a public polynomial-system mechanism for finding and certifying
the required information in \QRUOV{}.

\paragraph{Why vinegar-only retuning cannot help.}
The direct-forgery and key-recovery routes respond oppositely to $V$.  When
$V\leq m-\ell-1$, key recovery uses at most $m-\ell-1$ variables.
When $V\geq m-\ell$, the direct reduction and fixing the common scale use exactly
$m-\ell-1$.  These cases cover every integer $V$, so
\[
 r_{\mathrm{attack}}(V)\leq m-\ell-1.
\]
For \QRUOV{}, $\ell=3$, so the ceiling is $m-4$
(\cref{thm:security-scissors}).  Increasing only $V$ therefore cannot move the
two attacks past their crossing point.  A parameter revision must also
increase $m$, or change the construction so that one of the reductions no
longer applies.

\paragraph{Concrete estimates and executable evidence.}
The reductions above still leave large nonlinear systems.  To estimate the
work required to solve them, we adapt an exact symbolic equation solver to
the base field of characteristic $127$.  The solver starts from equations
whose roots are known, changes them gradually into the attack equations, and
tracks the roots through this change.  We also account for its polynomial
arithmetic over extension fields.  Under the specification's Boolean-gate
cost convention, the resulting estimates are
$\mathbf{2^{140.384}}$, $2^{191.027}$, and $2^{247.125}$ gates at Levels I,
III, and V, respectively.  These rows lie $2.616$, $15.973$, and $24.875$ bits
below their category targets (\cref{tab:combined-costs}).  The Level-I row is a
candidate break only within this cost model.  It depends on extrapolating a
degree bound to the exact attack equations and on an unimplemented fast
polynomial multiplication schedule.  The symbolic representation also needs
an impractical amount of memory (\cref{sec:concrete-solver,sec:security}).

On smaller examples, the direct-forgery implementation runs from public input
to a forged signature accepted by the verifier over two small fields.  This
checks that the attack steps work together, but it does not validate the
full-parameter cost estimate (\cref{sec:evidence-limitations}).  The algebraic
reductions and final verification are exact.  The success probabilities use
the stated random-key models, and the gate counts remain model-based
estimates.

\paragraph{Organization.}
\Cref{sec:related,sec:model} give the prior-work and specification context.
\Cref{sec:direct,sec:recovery} prove the two reductions and the fixed-output
bound.  The remaining sections give the root solver, concrete estimates,
evidence, limitations, and parameter consequences.

\section{Related work}\label{sec:related}

\paragraph{From Oil and Vinegar to \QRUOV{}.}
UOV was introduced after the balanced construction's hidden linear structure
was recovered; its vinegar imbalance must be chosen against both structural
attacks and an increasingly underdetermined public system
~\cite{KipnisShamir1998,KPG1999,BraekenWolfPreneel2005}.  Field lifting and
block-anti-circulant matrices compress the public key
~\cite{FieldLifting,BACUOV}, but the latter construction was broken through
its reducible quotient polynomial~\cite{BACUOVAttack}.  \QRUOV{} instead uses
an irreducible quotient and analyzes the resulting pull-back and lifting
attacks~\cite{QRUOVOriginal}.  We hold this construction fixed and vary only
its vinegar dimension.

\paragraph{Direct attacks on underdetermined quadratic systems.}
Thomae and Wolf reduce an underdetermined quadratic system to a smaller square
system before invoking a generic solver~\cite{ThomaeWolf2012}.  Furue,
Nakamura, and Takagi improve the range in which this reduction applies, and
Hashimoto divides the residual equations into smaller blocks
~\cite{FurueUnderdetermined2021,Hashimoto2023}.  Hashimoto's method is the
classical benchmark used by the QR-UOV team.  Its January 2026 response gives
costs of $158$, $217$, and $285$ bits for the recommended parameter sets
~\cite{QRUOVResponse2026}.  The resulting affine residual dimensions
$51,75,102$ are therefore prior art.  The later ``Just Guess'' algorithm
pushes more of the computation into exhaustive search.  It improves the
quantum estimates for QR-UOV, but it does not improve the team's classical
Hashimoto estimates~\cite{JustGuess,QRUOVResponse2026}.  Our direct attack
starts from the same dimensions.  We change the output coordinates so that
the remaining target equations are homogeneous, which means that solutions
are defined only up to a common nonzero scale.  Fixing that scale removes one
variable and gives systems in $50,74,101$ variables.

\paragraph{Recovering the hidden oil space.}
The Kipnis--Shamir attack searches directly for the secret linear structure
~\cite{KipnisShamir1998}.  Reconciliation expresses membership in the oil
space as quadratic equations.  Beullens' intersection attack instead compares
several transformed copies of the oil space and uses their intersections to
reduce the search
~\cite{DingEtAl2008,Beullens}.  P\'ebereau strengthens the endpoint: for the
practical UOV parameters he analyzes, one suitable oil vector is enough to
recover an equivalent key in polynomial time~\cite{Pebereau}.  We do not claim this
one-vector endpoint.  Our contribution is a public QR-UOV mechanism for
finding the required information.  A chosen extension-field direction exposes
a planted root $Gc$, the derivative test certifies when that root determines
the hidden graph, and exact signing checks certify the recovered equivalent
key (\cref{sec:recovery}).

\paragraph{QR-UOV-specific key recovery.}
The original QR-UOV analysis applies reconciliation and intersection after
pulling the block matrices back to the quotient field or lifting the public
map to an extension field~\cite{QRUOVOriginal}.  Furue and Ikematsu instead
encode the same structure as the problem of finding a low-rank matrix in a
public linear space
~\cite{FurueIkematsu}.  Suzuki et al. extend the admissible target ranks and
connect this method to the singular points used by other UOV attacks
~\cite{Suzuki}.  More recent methods derive higher-degree algebraic relations
using symmetric powers or truncated polynomial rings.  These relations
improve key recovery for some UOV parameter sets
~\cite{JinEtAl2026,FurueIkematsu2026}.  For the recommended QR-UOV sets, the
truncated-ring calculation gives no improvement because the characteristic
$127$ already exceeds the relevant solving degree, and the reported
odd-characteristic symmetric-algebra costs do not beat reconciliation.  These
works sharpen individual key-recovery costs at fixed parameters.  None gives a
bound that joins key recovery to direct forgery for every vinegar dimension.

\paragraph{Parameter tradeoffs and common UOV frameworks.}
P\'ebereau and the original QR-UOV parameter analysis already state the
qualitative tension between key recovery and forgery
~\cite{Pebereau,QRUOVOriginal}.  Li, Guo, and Ding describe QR-UOV, MAYO, and
SNOVA in one algebraic framework and show that QR-UOV key-recovery costs can
be analyzed as UOV over the extension field~\cite{LiGuoDing2025}.  We claim neither the
tradeoff nor the extension-field view.  Our result is the exact two-case bound:
every integer $V$ leaves either the direct or equivalent-key reduction with at
most $m-\ell-1$ variables, and the direct reduction reaches this bound only
after fixing the common scale (\cref{thm:security-scissors}).

\paragraph{Polynomial-system solving.}
We estimate root-finding costs with the exact symbolic algorithm of Safey El
Din and Schost~\cite{SafeySchost2018}.  It connects an easy system with known
roots to the target system and follows those roots as the connecting parameter
changes~\cite{Schost2003,GiustiLecerfSalvy2001}.  We change the easy starting
system and the linear form used to distinguish roots so that the algorithm
works in characteristic $127$.  The supplement checks the result against two
other exact solvers~\cite{vdHL,Gimenez2026}.  These solvers estimate the cost
of the systems produced by our attacks.  They do not produce either QR-UOV
reduction or the fixed-output theorem.

\paragraph{Chevalley--Warning.}
The direct reduction uses the classical fact that, when the sum of the degrees
is smaller than the number of variables, the number of common zeros over a
finite field is divisible by the characteristic~\cite{Chevalley1935,Warning1935}.
Applied to three quadrics in seven variables, it guarantees a nonzero extension
vector.  Exhaustive search over nonzero vectors up to scalar multiplication
makes the theorem constructive, so the preprocessing does not assume that a
common zero is available for free.  Our use of this classical theorem is
therefore algorithmic rather than a new finite-field existence result.

\section{From a recovered graph to signatures on any message}\label{sec:model}

The direct attack can check a proposed signature by evaluating the public map
$\widehat P$.  Key recovery has two additional obligations.  A recovered
graph must reproduce the same public map, and its oil system must let the
attacker solve for every verification target.  This section states the random
key models used only for success probabilities, then gives public checks for
both obligations.

\subsection{Parameters and expansion models}

The Round-2 parameter sets specialize \eqref{eq:parameter-stack} to $q=127$
and $\ell=3$.  Thus
\[
 K=\F_{q^\ell}=\F_{127^3},\qquad
 Q=|K|=2{,}048{,}383,
 \qquad \log_2 Q=20.966054\ldots.
\]
The three parameter sets are as follows.
\begin{center}
\begin{tabular}{@{}lrrrrrr@{}}
\toprule
Level & $v$ & $m$ & $V$ & $O$ & $N$ & $e=m-V$\\
\midrule
I&156&54&52&18&70&2\\
III&228&78&76&26&102&2\\
V&306&105&102&35&137&3\\
\bottomrule
\end{tabular}
\end{center}
We write $D=2^V$ for B\'ezout's general upper bound on the number of isolated
roots of $V$ quadratic equations in $V$ variables.  We write $L_s=\F_{Q^s}$ for an
extension field used by the equation solver.  All concrete theorems below
concern these three sets, for which $m<q$.  The matrices
$A_i,B_i,E_i,C_i$, the graph $G$, and the oil space $\cO$ retain the meanings
in \eqref{eq:central} to \eqref{eq:isotropy}.

For the probability calculations, we replace the seeded expansion by the
following explicit distribution on expanded keys.  The algebraic attacks and
all final verification checks do not depend on this replacement.

\begin{definition}[Random expanded-key model]\label{def:ideal}
Sample the symmetric blocks $A_1,\ldots,A_m$ independently and uniformly from
$\Sym_V(K)$.  Conditional on $A=(A_i)_i$, sample the secret graph
$G\in K^{V\times O}$ from any distribution.  Conditional on $(A,G)$, sample
$E_1,\ldots,E_m$ independently and uniformly from $K^{V\times O}$, and set
each $C_i$ by \eqref{eq:graphid}.  No independence between $G$ and $A$ is
assumed.
\end{definition}

\begin{definition}[Independent-graph model]\label{def:direct-ideal}
The independent-graph model is Definition~\ref{def:ideal} with $G$ independent of $A$.
The marginal distribution of $G$ may otherwise be arbitrary.
\end{definition}

\begin{proposition}[Central blocks remain uniform]\label{prop:ideal-factorization}
In Definition~\ref{def:ideal}, the blocks $B_i=E_i-A_iG$ are independent and uniform in
$K^{V\times O}$, and $B=(B_i)_i$ is independent of $(A,G)$.  Hence the pairs
$(A_i,B_i)_{1\le i\le m}$ are independent and uniform on
$\Sym_V(K)\times K^{V\times O}$, although $G$ may remain correlated with $A$.
\end{proposition}

\begin{proof}
For fixed $(A,G)$, each map $E_i\mapsto E_i-A_iG$ is a translation of
$K^{V\times O}$.  It therefore preserves the uniform law and makes that law
independent of $(A,G)$.  Independence of the $A_i$ and $E_i$ gives the stated
product distribution.
\end{proof}

The equivalent-key analysis uses the first model.  Its failure events depend
on the uniform central matrices and on fresh public attack randomness, but not
on independence between $G$ and $A$.  The direct reduction is valid for every
key.  Its average success probability, however, requires $G$ to be independent
of the central coefficients and therefore uses the second model.  The
supplement explains when idealized random functions or ideal ciphers justify
these distributions and why the argument does not automatically transfer to
the fixed SHAKE or AES functions (\cref{app:seeded-transfer}).  These models
affect only the probability of finding a candidate.  Every returned key and
signature is checked against the original public map.

\subsection{From a recovered graph to a signing key}

The graph identities are useful because they can be checked using only public
matrices.  Passing these checks reconstructs a central UOV map for the same
public key.

\begin{proposition}[A verified graph gives an exact UOV decomposition]\label{prop:verified-key}
Let $G'\in K^{V\times O}$ satisfy every identity \eqref{eq:graphid}, and put
\[
 B_i'=E_i-A_iG',\qquad
 F_i'=\begin{pmatrix}A_i&B_i'\\{B_i'}^{\trans}&0\end{pmatrix},\qquad
 S_{G'}=\begin{pmatrix}I_V&G'\\0&I_O\end{pmatrix}.
\]
Then $P_i=S_{G'}^{\trans}F_i'S_{G'}$ for every $i$.  The specification's
invertible base-field conversion therefore gives matrices
$\widehat P_i=\widehat S_{G'}^{\trans}\widehat F_i'\widehat S_{G'}$ over
$\F_q$ for the unchanged verification map~\cite[Sec.~4.1]{QRUOVSpec}.
\end{proposition}

\begin{proof}
Direct block multiplication gives the extension-field factorization because
the lower-right block is zero exactly when \eqref{eq:graphid} holds.  Applying
the specification's inverse pull-back gives the base-field factorization.  The
full matrix calculation appears in \cref{app:expanded-certificates}.
\end{proof}

\begin{proposition}[Public signing witness]\label{prop:signing-witness}
If $\det\mathsf L_{\widehat F}(a)\ne0$, then
$(\widehat F,\widehat S,a)$ is a publicly checkable signing certificate.  For
every target $y\in\F_q^m$, the unique solution to
\begin{equation}\label{eq:universal-inversion}
 o=\mathsf L_{\widehat F}(a)^{-1}\bigl(y-\widehat F(a,0)\bigr),
 \qquad
 x=\widehat S^{-1}(a,o)
\end{equation}
satisfies $\widehat P(x)=y$.
\end{proposition}

\begin{proof}
Substitution gives $\widehat F(a,o)=y$ and then
$\widehat P(x)=\widehat F(\widehat Sx)=y$.  The forged signatures need not
match the honest signing distribution.
\end{proof}

\begin{corollary}[Forging any chosen message without signing queries]\label{cor:universal-forgery}
A verified certificate from Proposition~\ref{prop:signing-witness} lets an attacker forge
an accepted \QRUOV{} signature on every chosen message without a signing
query.  Given a message $M$, choose any allowed salt $r$, form $y(M,r)$ by
\eqref{eq:verification-target}, invert it with
Proposition~\ref{prop:signing-witness}, and output $(r,x)$.
\end{corollary}

\begin{proof}
The certificate satisfies the verifier equation
\eqref{eq:verification-equation} for every target.  It neither recovers nor
uses the secret seed.
\end{proof}

\subsection{Finding vinegar values with an invertible oil system}

Fix nonzero $\xi\in\F_q^\ell$.  For each quotient-ring vinegar block $j$, let
$a_j$ place $\xi$ in block $j$ and zero elsewhere, and set
\[
 M_j:=\mathsf L_{\widehat F}(a_j)\in\F_q^{m\times m}.
\]
Write
\[
 \pi_m:=\prod_{k=1}^m(1-q^{-k}),
 \qquad \sigma_m:=1-\pi_m.
\]
$\pi_m$ is the probability that a uniform $m\times m$ matrix over $\F_q$ is
invertible.  We first test vinegar assignments supported on one quotient-ring
block.  This makes the search both public and inexpensive.

\begin{proposition}[Searching basis vinegar assignments]\label{prop:signing-search}
For the planted decomposition in Definition~\ref{def:ideal}, the matrices
$M_1,\ldots,M_V$ are independent and uniform.  The following searches then
give public signing certificates.
\begin{enumerate}
\item Testing basis assignments in order uses fewer than
$\pi_m^{-1}<1.008$ tests on average.  Testing all $V$ fails with probability
$\sigma_m^V$, whose base-two logarithms at Levels I, III, and V are
$-362.823$, $-530.280$, and $-711.692$.

\item Testing only the first $22$, $31$, and $40$ assignments at the three
levels fails with base-two log probabilities
\begin{equation}\label{eq:compact-sign-logs}
 -153.502,\qquad-216.298,\qquad-279.095.
\end{equation}

\item Pairing the basis matrices as $(M_1,M_2),(M_3,M_4),\ldots$ gives a
deterministic fallback.  Let $\cE_{\rm sign}$ be the event that every chosen
pair span contains only singular matrices.  Its probability satisfies
\begin{equation}\label{eq:eta-sign}
 \Pr[\cE_{\rm sign}]\le\eta_{\rm sign}.
\end{equation}
Its base-two logarithms at Levels I, III, and V are below
\[
 -544.820,\qquad-796.276,\qquad-1068.687.
\]
Outside this event, testing every $M_j$ and then
$A+tB$ for each chosen pair and $t\in\F_q^\times$ succeeds after at most
$V+\lfloor V/2\rfloor(q-1)$ tests.
\end{enumerate}

For any exact candidate, define
\[
 \Delta(c_1,\ldots,c_V):=\det\!\left(\sum_{j=1}^V c_jM_j\right).
\]
The coefficients $(c_1,\ldots,c_V)$ correspond to the vinegar assignment
$a=(c_1\xi,\ldots,c_V\xi)$.  If $\Delta$ is nonzero, a uniform assignment of
this form succeeds
with probability at least $1-m/q$.  The expected numbers of assignments are
at most $1.740$, $2.592$, and $5.773$ at the three levels.  This last claim is
key by key, and every successful determinant test is an exact public
certificate even when the recovered graph differs from the original secret
graph.
\end{proposition}

\begin{proof}
The specification's regular representation maps each independent uniform
entry of the central block $B_i$ bijectively to one row block of $M_j$.  Thus
the $M_j$ are independent and uniform.  The remaining statements follow from
the exact rank distribution of a uniform matrix, a count of singular
two-dimensional spans, and Schwartz--Zippel.  The formulas and calculations
appear in \cref{app:expanded-certificates}.
\end{proof}

\section{\texorpdfstring{Direct forgery in $m-4$ variables}{Direct forgery in m-4 variables}}\label{sec:direct}

The direct attack starts from the verifier equation
\eqref{eq:verification-equation} over $\F_q$, with $q=127$ and
$n=v+m$.  Given a nonzero target $y$, it constructs an $(m-3)$-dimensional
public subspace on which three zero-target equations vanish identically.  On
this subspace, the remaining zero-target equations are homogeneous.  Their
solutions are therefore defined up to multiplication by a nonzero scalar.
Fixing this scale leaves $m-4$ equations in $m-4$ variables.  The final
nonzero-target equation recovers the omitted scale.  The resulting systems
have dimensions
\[
 50,\qquad74,\qquad101
\]
for Levels I, III, and V.

\begin{remark}[Relation to Hashimoto and the specification]\label{rem:direct-prior-art}
The affine reduction to $m-3$ variables is prior art.  Hashimoto's
underdetermined quadratic-system method and the QR-UOV specification already analyze these
residual dimensions~\cite{Hashimoto2023,QRUOVSpec}.  Our additional reduction
uses full target normalization and solves the homogeneous zero-target
equations modulo scaling.  This removes one more variable, giving systems of
dimension $m-4$.  We also give a finite construction of the required
subspace, bound its failure probability for every fixed key, prove a lower
bound on the probability of a root with invertible Jacobian, and estimate the
cost over the base field of characteristic $127$.
\end{remark}

The algebraic reduction below is key by key.  Its probability calculation uses
the independent-graph model of Definition~\ref{def:direct-ideal}: the
hidden graph $G$ is independent of the independent uniform central coefficient
pairs $(A_i,B_i)$.  This independence is needed only for the direct attack's
average success analysis.  The equivalent-key results continue to allow
arbitrary correlation between $G$ and $A$.  In idealized models where SHAKE is
a random function or AES is an ideal cipher, independently sampled and
distinct public and secret seeds give the required independence.
\Cref{app:seeded-transfer} accounts for seed collisions and finite output
lengths.  We do not claim that fixed SHAKE or AES automatically has the same
distribution.

\subsection{From the target to a projective system}

If $y=0$, then $x=0$ is already a valid preimage, so only the nonzero case
needs analysis.  Let $y\in\F_q^m$ be nonzero.  Choose a public matrix
$T_y\in\operatorname{GL}_m(\F_q)$ whose first $m-1$ rows span $y^\perp$ and
whose last row has inner product one with $y$.  Then
\begin{equation}\label{eq:target-normalization}
 T_yy=(0,\ldots,0,1)^{\mathsf T}.
\end{equation}
Write $Q_1,\ldots,Q_m$ for the corresponding output recombinations of the
public forms.  Since $T_y$ is fixed by $y$ and invertible, independent uniform
public forms remain independent and uniform in the random expanded-key model.

Only $Q_1,Q_2,Q_3$ are used to construct the attack subspace.  Choose a
uniformly random seven-dimensional subspace $S_0\subset\F_q^n$ and restrict
these three forms to $S_0$.  The attacker then tests all nonzero vectors in
$S_0$ up to scalar multiplication, namely the points of
$\PP(S_0)\simeq\PP^6(\F_q)$.  Chevalley--Warning guarantees a nonzero common
zero because three quadrics have total degree six in seven variables.  Denote
the first such point by $w$.

The hidden oil space is not known to the attacker, so the algorithm cannot
test whether $w$ lies outside it.  We call such a point \emph{off oil}.  The
random seed subspace makes this test unnecessary; the next proposition bounds
the failure probability for every fixed oil space.

\begin{proposition}[A random seed misses every fixed oil space]\label{prop:off-oil-seed}
Fix an arbitrary $m$-dimensional subspace $\widehat{\cO}\subset\F_q^n$ and
choose a uniform seven-dimensional subspace $S_0\subset\F_q^n$.  Then
\begin{equation}\label{eq:seed-oil-bound}
 \Pr[S_0\cap\widehat{\cO}\ne\{0\}]
 \le
 \epsilon_{\rm seed}
 :=\frac{q^m-1}{q-1}\,
   \frac{q^7-1}{q^n-1}.
\end{equation}
Consequently, except with probability at most $\epsilon_{\rm seed}$ over
the attack's public randomness, every nonzero common zero found inside $S_0$ is off oil.
For the three official profiles,
\begin{equation}\label{eq:seed-oil-numbers}
 \log_2\epsilon_{\rm seed}
 \le -1048.291,\qquad-1551.477,\qquad-2096.594.
\end{equation}
The bound holds for every fixed hidden oil space; it does not average over key
generation.
\end{proposition}

\begin{proof}
There are $(q^m-1)/(q-1)$ projective oil points.  For any fixed projective
point $[u]$, a uniform seven-dimensional subspace contains $[u]$ with
probability $(q^7-1)/(q^n-1)$.  A union bound proves
\eqref{eq:seed-oil-bound}.  The numerical values substitute the three official
pairs $(n,m)$.
\end{proof}

\paragraph{Constructing the isotropic subspace.}

For a quadratic form $Q$ over a field of odd characteristic, write
\[
 B_Q(u,z)=\frac{Q(u+z)-Q(u)-Q(z)}2
\]
for its polar bilinear form.  A subspace is \emph{totally isotropic} for $Q$ when $Q$ vanishes identically on it.

\begin{theorem}[Building a common isotropic subspace]\label{thm:cw-extension}
Let $Q_1,\ldots,Q_s$ be homogeneous quadrics on $\F_q^n$, and let $0\ne w$ satisfy $Q_i(w)=0$ for every $i$.  If
\begin{equation}\label{eq:cw-feasibility}
 n\ge(s+1)r+s,
\end{equation}
then a deterministic finite algorithm constructs an $r$-dimensional subspace $U$ containing $w$ on which every $Q_i$ vanishes identically.

At each extension step the algorithm uses linear algebra and examines at most
\begin{equation}\label{eq:cw-projective-count}
 \left|\PP^{2s}(\F_q)\right|
 =\frac{q^{2s+1}-1}{q-1}
\end{equation}
projective points.
\end{theorem}

\begin{proof}
At each step, use linear equations to restrict to vectors $z$ satisfying
$B_{Q_i}(u,z)=0$ for every current basis vector $u$ and every selected form
$Q_i$.  Condition \eqref{eq:cw-feasibility} leaves at least $2s+1$
quotient dimensions, where Chevalley--Warning supplies a nonzero common zero
of the induced $s$ quadrics~\cite{Chevalley1935,Warning1935}.  Lifting this
zero extends the isotropic subspace by one dimension.  The complete dimension
count appears in \cref{app:direct-proofs}.
\end{proof}

\begin{corollary}[Affine reduction after target normalization]\label{cor:generic-target-slicing}
Let $P:\F_q^n\to\F_q^m$ be any homogeneous quadratic map, let $y\ne0$, and
choose $1\le s<m$.  After an invertible output recombination sends $y$ to the
last coordinate, a nonzero common zero of any $s$ of the first $m-1$
transformed quadrics can be extended to an $(m-s)$-dimensional common totally
isotropic subspace whenever
\begin{equation}\label{eq:generic-square-condition}
 n\ge(s+1)m-s^2.
\end{equation}
On that subspace, the affine problem has $m-s$ equations in $m-s$ variables.
\end{corollary}

\begin{proof}
Apply \cref{thm:cw-extension} with $r=m-s$.  Its condition becomes
$n\ge(s+1)(m-s)+s=(s+1)m-s^2$.  The selected $s$ forms vanish identically on
the constructed subspace and have zero target, leaving the stated affine
residual.
\end{proof}

The affine formulation still counts a common scalar that does not change any
homogeneous zero-target equation.  A projective point $[z]$ records a nonzero
vector $z$ only up to multiplication by a nonzero scalar.  The next corollary
uses this viewpoint to fix the scale and remove one variable exactly.

\begin{corollary}[Projective reduction after target normalization]\label{cor:projective-target-slicing}
Assume that $q$ is odd.  Under the hypotheses of
Corollary~\ref{cor:generic-target-slicing}, put $r=m-s$ and let
$U$ be the constructed $r$-dimensional subspace.  Normalize the target to
$(0,\ldots,0,1)$ and write
\[
 H_1,\ldots,H_{r-1},H_0
\]
for the restrictions to $U$ of the remaining transformed forms, with $H_0$
corresponding to the nonzero target coordinate.  Then:
\begin{enumerate}
\item affine solutions of
\begin{equation}\label{eq:projective-affine-pair}
 H_1(z)=\cdots=H_{r-1}(z)=0,
 \qquad H_0(z)=1
\end{equation}
modulo the pair $z\sim-z$ are in bijection with projective points
$[z]\in\PP(U)$ satisfying
\begin{equation}\label{eq:projective-zero-system}
 H_1(z)=\cdots=H_{r-1}(z)=0
\end{equation}
and $H_0(z)$ a nonzero square in $\F_q$;
\item a solution of \eqref{eq:projective-affine-pair} has invertible affine
Jacobian if and only if the Jacobian of
\eqref{eq:projective-zero-system} has full rank at its projective point; and
\item after choosing a uniform nonzero linear form $L\in U^*$, the chart
$L(z)=1$ turns \eqref{eq:projective-zero-system} into $r-1$ quadrics in
$r-1$ affine variables.  Any fixed projective root lies in this chart with
probability
\[
 1-\frac{q^{r-1}-1}{q^r-1}>1-\frac1q.
\]
\end{enumerate}
Thus the system has $m-s-1$ equations in $m-s-1$ variables.  B\'ezout's
theorem bounds its number of isolated roots by $2^{m-s-1}$.
\end{corollary}

\begin{proof}
Homogeneity multiplies every quadratic value by $\lambda^2$ under
$z\mapsto\lambda z$.  Thus $H_0$ can be scaled to one exactly when its value
is a nonzero square.  Euler's identity shows that the affine Jacobian is
invertible exactly when the projective Jacobian has full rank after removing
the common scaling direction.  Counting the
linear forms that vanish at a fixed projective point gives the chart
probability.  The full Jacobian argument appears in
\cref{app:direct-proofs}.
\end{proof}

For the QR-UOV parameters, set $s=3$ and $r=m-3$.  The subspace condition is
\begin{equation}\label{eq:official-cw-condition}
 n\ge4m-9.
\end{equation}
The three official profiles give
\[
 (n,4m-9)=(210,207),\quad(306,303),\quad(411,411).
\]
The resulting projective solver dimensions are
\begin{equation}\label{eq:official-projective-cores}
 k=r-1=m-4\in\{50,74,101\}.
\end{equation}
The seed root and every extension step use finite exhaustive search, so this
reduction needs no existence oracle.  Even after overcharging the required
linear algebra and search, preprocessing costs at most $2^{64.501}$ Boolean
gates across the three profiles.  The complete cost calculation appears in
\cref{app:direct-proofs}.

\subsection{Why the reduced system has a usable root}

The structured degree-three conversion does not produce uniformly random
quadratic forms over $\F_q$.  We therefore prove the exact randomness property
needed for the root count instead of treating the reduced equations as random.
Let $E=K^{V+O}$, viewed over $\F_q$, and let
$\widehat{\cO}\subset E$ be the hidden base-field oil space.  In central
coordinates, each form in the random expanded-key model has the form
\begin{equation}\label{eq:trace-pullback}
 Q(v,o)=\Tr_{K/\F_q}\!\left(\gamma
 (v^{\mathsf T}Av+2v^{\mathsf T}Bo)\right)
\end{equation}
for fixed nonzero $\gamma\in K$ and uniform
$A\in\Sym_V(K)$ and $B\in K^{V\times O}$.

\begin{lemma}[Randomness of values and first derivatives]\label{lem:pullback-one-jet}
Let $x\in E\setminus\widehat{\cO}$, and let $U\subseteq E$ contain $x$.  For a
uniform form $Q$ from \eqref{eq:trace-pullback}, $Q(x)$ is uniform in $\F_q$.
Conditioned on $Q(x)=a$, the derivative $dQ_x|_U$ is uniform among the linear
forms $\lambda\in U^*$ satisfying $\lambda(x)=2a$.  Moreover, evaluation at
two off-oil points is linearly dependent only when the points differ by a
nonzero scalar in $\F_q$.
\end{lemma}

\begin{proof}
Nondegeneracy of the trace pairing lets the central bilinear form prescribe
every value and derivative allowed by Euler's identity.  Comparing the
rank-one coefficient matrices at two points gives the last claim.  The full
calculation of all possible values and derivatives, including the conversion
from the degree-three field, appears in
\cref{app:direct-proofs}.
\end{proof}

Let $U$ be the subspace from \cref{thm:cw-extension}, and choose an $\F_q$-basis
to identify $U$ with $\F_q^r$.  Under the normalization
\eqref{eq:target-normalization}, write
\[
 H_j=Q_{3+j}|_U\quad(1\le j\le r-1),
 \qquad H_0=Q_m|_U.
\]
The remaining affine system is
\begin{equation}\label{eq:direct-residual}
 H_1(z)=\cdots=H_{r-1}(z)=0,
 \qquad H_0(z)=1,
 \qquad r=m-3.
\end{equation}
The projective solver uses only the first $r-1$ equations; $H_0$ filters and
scales their projective roots.

The next statement gives an explicit success probability rather than assuming
that a system which looks random has a root.  Let
\begin{equation}\label{eq:pi-r-direct}
 \pi_{r-1}(q)=\prod_{j=1}^{r-1}(1-q^{-j}).
\end{equation}

\begin{proposition}[Probability of an affine root with invertible Jacobian]\label{prop:direct-root-probability}
Condition on the selected forms, the public randomness used to construct $U$, and
an off-oil seed point.  Put
\[
 d=\dim_{\F_q}(U\cap\widehat{\cO}),
 \qquad
 \mu=1-q^{d-r}.
\]
Then $d\le r-1$.  Over the independent remaining forms in the random
expanded-key model,
the probability that \eqref{eq:direct-residual} has an $\F_q$-rational root
with invertible $r\times r$ Jacobian is at least
\begin{equation}\label{eq:direct-root-lower-bound}
 \frac{\mu\,\pi_{r-1}(q)^2}
 {2+\mu-(q-1)q^{-r}}.
\end{equation}
For $q=127$ and each of $r=51,75,102$, this is at least
\begin{equation}\label{eq:direct-root-number}
 p_*=0.326336998767509\ldots.
\end{equation}
\end{proposition}

\begin{proof}
The off-oil seed point gives $d\le r-1$.  Let $M$ count all off-oil roots of
\eqref{eq:direct-residual}.  Lemma~\ref{lem:pullback-one-jet} gives
$\mathbb E[M]=\mu$.  Pairing proportional and nonproportional points gives
\begin{equation}\label{eq:root-second-moment}
 \mathbb E[M^2]
 =2\mu+\mu^2-(q-1)\mu q^{-r}.
\end{equation}
At a fixed root, the same lemma makes the Jacobian invertible with probability
$\pi_{r-1}(q)$.  If $N$ counts these roots, then
$\mathbb E[N]=\mu\pi_{r-1}(q)$ and $N^2\le M^2$.  The second-moment inequality
now gives \eqref{eq:direct-root-lower-bound}; using
$\mu\ge1-q^{-1}$ gives \eqref{eq:direct-root-number}.  The complete pair count
appears in \cref{app:direct-proofs}.
\end{proof}

\begin{corollary}[Probability of a usable projective root]\label{cor:projective-root-probability}
Under the same conditioning, the probability that the projective system
\[
 H_1=\cdots=H_{r-1}=0\quad\text{in }\PP^{r-1}(\F_q)
\]
has a point $[z]$ with full-rank projective Jacobian for which $H_0(z)$ is a nonzero square is at
least the bound in \eqref{eq:direct-root-lower-bound}, and hence at least
$p_*$ for every official profile.
\end{corollary}

\begin{proof}
By Corollary~\ref{cor:projective-target-slicing}, every affine root of
\eqref{eq:direct-residual} with invertible Jacobian maps to a projective root
with full-rank Jacobian and the
required nonzero-square value;
the two affine roots $z$ and $-z$ map to the same projective point.  Existence
of an affine root therefore implies the displayed projective event.
\end{proof}

The factor $1/p_*=3.0643169600\ldots$ adds $1.615566$ bits to the expected
cost.  A random chart adds less than
$\log_2(127/126)=0.011405$ bits.  The root bound is proved from the QR-UOV
distribution by a second-moment argument.  The probability that the initial
subspace avoids the oil space holds for every fixed key.  The later root
probability averages over the independent-graph model.  We do not claim that
repeated attacks on every fixed seeded key behave independently.

\subsection{The verified attack}

A projective chart fixes the common scale by choosing a nonzero linear form
$L$ and requiring $L=1$.  With this terminology, the complete direct attack
is as follows.  The root solver also applies a random linear form to label its
candidate roots.  We call this form a \emph{separator} when the roots that the
algorithm must distinguish receive different labels.

\begin{theorem}[Verified QR-UOV direct forgery]\label{thm:verified-direct-forgery}
Work in the independent-graph random expanded-key model and the official degree-three
pull-back.  For any nonzero target $y$, the following public algorithm has
one-sided error: it may fail to find a forgery, but it never returns an invalid
one.
\begin{enumerate}
\item Normalize the target to $(0,\ldots,0,1)$ as in
\eqref{eq:target-normalization}.
\item Choose a uniform seven-dimensional seed subspace, enumerate its
projective points, and take a nonzero common zero of the first three forms.
\item Use \cref{thm:cw-extension} to construct an $(m-3)$-dimensional common
totally isotropic subspace for those forms.
\item Choose a uniform nonzero chart form $L$ on that subspace.  Substitute
$L=1$ in the $m-4$ remaining zero-target quadrics, and apply
\cref{thm:finite-char-nonsingular} to the resulting square system in
\[
 k=m-4
\]
variables.
\item For each $\F_q$-rational root returned by the solver, recover its projective point,
test whether the last quadratic value is a nonzero square, scale to make that
value one, reverse the public transformations, and evaluate the unchanged
public map.
\end{enumerate}
Every returned vector is a valid signature preimage.  Over the random
expanded-key model and the attack's public randomness, one complete invocation
reaches a reduced system with a root over $\F_q$ whose Jacobian is invertible
and whose last quadratic value is a nonzero square with probability at least
\begin{equation}\label{eq:direct-total-success}
 (1-\epsilon_{\rm seed})p_*
 \left(1-\frac1q\right)
 \left(1-\frac{2^{2k}}{|k'|}\right),
 \qquad k=m-4,
\end{equation}
where $k'/\F_q$ is the solver working extension.  Taking
\[
 [k':\F_q]=15,\qquad22,\qquad30
\]
at Levels I, III, and V gives the finite costs reported in
\cref{sec:security}.  We multiply the cost of one invocation by the inverse of
\eqref{eq:direct-total-success}.  This is a conservative expected-attempt
factor, not a claim that retries on every fixed key are independent.
\end{theorem}

\begin{proof}
The first $m-1$ target coordinates are zero.  By
Proposition~\ref{prop:off-oil-seed}, the seed subspace is disjoint from the hidden oil
space with probability at least $1-\epsilon_{\rm seed}$ for every fixed key.
Chevalley--Warning then supplies an off-oil seed root, and
\cref{thm:cw-extension} constructs the advertised subspace.  Conditional on
the selected first three forms, the graph, the public randomness, and this off-oil
seed point, the remaining transformed forms are still independent uniform
pull-back forms.  Corollary~\ref{cor:projective-root-probability} supplies a
usable projective root with full-rank Jacobian with probability at least
$p_*$.  A random nonzero chart contains a fixed such root with probability
greater than $1-1/q$.  The solver of \cref{sec:finite-char-solver} includes
every root in the chart with the separator probability in
\eqref{eq:direct-total-success}.  Square testing and
scaling recover the associated affine solutions exactly.  Re-evaluation of
all $m$ original public forms rejects every invalid candidate,
so the algorithm cannot output an invalid signature.
\end{proof}

\subsection{Why changing only the vinegar count cannot help}

Combining the direct reduction with equivalent-key recovery gives
a bound for every vinegar dimension.  Let the quotient degree be $\ell$, so that
\[
 n=\ell V+m.
\]
Applying Corollary~\ref{cor:generic-target-slicing} with a general value of $s$ gives
the required $(m-s)$-dimensional common-isotropic subspace whenever
\begin{equation}\label{eq:scissors-feasibility}
 \ell V\ge s(m-s).
\end{equation}
By Corollary~\ref{cor:projective-target-slicing}, the corresponding nonzero-target
square system has $m-s-1$ variables.

\begin{theorem}[Fixed-output bound across all vinegar dimensions]\label{thm:security-scissors}
Work over a field of odd characteristic.  Fix $m$ and quotient degree $\ell$.  For every nonzero target, within the two
attack families analyzed in this paper a square system in at most
$m-\ell-1$ variables is available for every integer vinegar count $V$:
\begin{equation}\label{eq:scissors-bound}
 r_{\rm attack}(V)\le m-\ell-1.
\end{equation}
More precisely:
\begin{enumerate}
\item if $V\le m-\ell-1$, equivalent-key recovery uses a square
system in $V$ variables;
\item if $V\ge m-\ell$, the direct reduction with $s=\ell$, followed by fixing
the common scale, uses a square system in $m-\ell-1$ variables.
\end{enumerate}
These cases cover every integer $V$.  For QR-UOV, $\ell=3$, so the ceiling is
$m-4$.
\end{theorem}

\begin{proof}
In the first case, $V<m$, so the $m$ key-recovery equations admit a
$V$-equation square recombination and the planted-root attack of
\cref{sec:recovery} has at most $V\le m-\ell-1$ variables.  In the second case,
\[
 \ell V\ge\ell(m-\ell),
\]
which is \eqref{eq:scissors-feasibility} with $s=\ell$.  The affine direct
system has dimension $m-\ell$ and
Corollary~\ref{cor:projective-target-slicing} removes the common-scale
coordinate, leaving $m-\ell-1$ variables.  Since consecutive integers
$m-\ell-1$ and $m-\ell$ meet the
two inequalities, no case is omitted.
\end{proof}

The theorem bounds the number of variables in the nonlinear system.  It does
not claim that running time depends only on this number.  For the official degree-three sets,
Proposition~\ref{prop:off-oil-seed} and
Corollary~\ref{cor:projective-root-probability} supply the success
factors charged in the concrete cost calculation.  Its design implication is exact:
once $V$ reaches $m-\ell$, increasing $V$ while holding $m$ fixed no longer
enlarges the projective direct-forgery system.  Conversely, decreasing $V$ into
the other case makes key recovery no larger than the same ceiling.
Moving both attack lines therefore requires changing $m$ as well as $V$, or
modifying the construction so that one reduction no longer applies.

The direct construction and the complete attack have both
been executed on reduced parameters.  The exact traces and public replay
checks appear in \cref{app:direct-proofs,app:reduced-implementation}; these are
correctness tests, not full-parameter benchmarks.

\section{Equivalent-key recovery from a planted root}\label{sec:recovery}

The second attack uses a root guaranteed by the hidden key structure.  Such a
guaranteed but unknown root is often called a \emph{planted root}.  Fixing an
oil vector $c\in K^O$ turns the hidden graph point
$(-Gc,c)\in\cO$ into the planted root $Gc$ of public quadratic equations in
$V$ variables.  If the matrix of first derivatives has full rank at that root,
public linear algebra recovers every column of $G$.  The attack then checks
\eqref{eq:graphid}, applies the pull-back \eqref{eq:official-pullback}, verifies
the factorization \eqref{eq:base-factorization}, and finds a signing witness
for \eqref{eq:universal-inversion}.  The end of this section records a second
way to recover the same oil space from the derivatives at a projective oil
point.

\subsection{Public equations with a hidden root}

Fix $c\in K^O$.  Substitute the public oil coordinate $c$ and the negated
vinegar coordinate $-x$ into each public form:
\begin{equation}\label{eq:direction-system}
 q_{i,c}(x):=P_i(-x,c)
 =x^{\mathsf T}A_ix-2x^{\mathsf T}E_ic+c^{\mathsf T}C_ic,
 \qquad x\in K^V.
\end{equation}
The equations $q_{1,c}=\cdots=q_{m,c}=0$ are public.  Because the choice of
$c$ fixes one direction in the oil coordinates, we call these the
\emph{fixed-direction equations}.  Define their public derivative matrix by
\begin{equation}\label{eq:direction-matrix}
 \mathcal D_c(x):=
 \begin{pmatrix}
 (E_1c-A_1x)^{\mathsf T}\\[-1mm]
 \vdots\\[-1mm]
 (E_mc-A_mx)^{\mathsf T}
 \end{pmatrix}.
\end{equation}
Recall that $B_i=E_i-A_iG$, and let $\mathcal B(c)$ have $i$th row
$(B_ic)^{\mathsf T}$.

\begin{proposition}[A planted root for every direction]\label{prop:graph-direction-identity}
For every $c\in K^O$, the point $x_c=Gc$ is a common zero of
\eqref{eq:direction-system}.  More precisely, for every $z\in K^V$,
\begin{equation}\label{eq:direction-translate}
 q_{i,c}(Gc+z)=z^{\mathsf T}A_iz-2(B_ic)^{\mathsf T}z.
\end{equation}
Consequently,
\[
 \mathcal D_c(x_c)=\mathcal B(c),
 \qquad
 J_q(x_c)=-2\mathcal B(c).
\]
\end{proposition}

\begin{proof}
Expand \eqref{eq:direction-system} at $x=Gc+z$.  The graph identity
\eqref{eq:graphid} makes the constant term zero, while the linear term is
$-2(B_ic)^{\mathsf T}z$.  This proves \eqref{eq:direction-translate};
differentiation gives the two matrix identities.
\end{proof}

\subsection{A root with full-rank derivative determines the graph}

A root $x$ for one oil vector $c$ gives only one graph value.  When the
derivative matrix has full column rank, the same public equations determine
the graph value on every other oil vector and therefore determine all of $G$.

\begin{proposition}[Exact graph completion]\label{prop:graph-direction-completion}
Let $x$ be a common zero of all equations $q_{i,c}$ and suppose
$\rank\mathcal D_c(x)=V$.  For $d\in K^O$, define
\begin{equation}\label{eq:direction-rhs}
 b_{c,d}(x)_i:=c^{\mathsf T}C_id-x^{\mathsf T}E_id
 \qquad(1\le i\le m).
\end{equation}
Then the public system
\begin{equation}\label{eq:direction-completion-system}
 \mathcal D_c(x)y=b_{c,d}(x)
\end{equation}
has at most one solution $y\in K^V$.  If $x=Gc$, its solution is $y=Gd$.
Solving \eqref{eq:direction-completion-system} for the $O$ coordinate vectors
$d=e_j$ therefore recovers $G$ exactly.  Checking every identity
\eqref{eq:graphid} then certifies a valid expanded UOV decomposition.
\end{proposition}

\begin{proof}
Full column rank gives uniqueness.  For $x=Gc$, multiply
\eqref{eq:graphid} by $c^{\mathsf T}$ on the left and $d$ on the right.  The
result is
\[
 (E_ic-A_iGc)^{\mathsf T}Gd
 =c^{\mathsf T}C_id-(Gc)^{\mathsf T}E_id,
\]
which is row $i$ of \eqref{eq:direction-completion-system}.  Hence $Gd$ is the
unique solution.  The verified-key conclusion follows from
Proposition~\ref{prop:verified-key}.
\end{proof}

\begin{remark}[A compact attack certificate]\label{rem:graph-certificate}
A successful candidate is certified by $(c,x,G)$ and one full-rank minor of
$\mathcal D_c(x)$.  A verifier checks the fixed-direction equations, the graph
completion systems, all graph identities, the inverse pull-back, and a signing
witness.  The certificate makes no claim about other roots or solver
completeness.
\end{remark}

\subsection{How often the derivative has full rank}

We now bound the probability that the planted root passes the full-rank test.
The following quantity is the exact full-rank probability for a uniform
$m\times V$ matrix over $K$.  Let
\begin{equation}\label{eq:rho-mv}
 \rho_{m,V}(Q):=\prod_{a=0}^{V-1}(1-Q^{a-m}),
 \qquad
 \tau_{m,V}(Q):=1-\rho_{m,V}(Q).
\end{equation}
Thus $\rho_{m,V}(Q)$ is the probability that a uniform $m\times V$ matrix over
$K$ has full column rank.

\begin{proposition}[Rank of the planted derivative]\label{prop:graph-direction-density}
For a fixed expanded key, suppose some maximal minor of $\mathcal B(c)$ is a
nonzero polynomial in $c$.  A fresh uniform $c\leftarrow K^O$ then satisfies
\begin{equation}\label{eq:random-direction-success}
 \Pr[\rank\mathcal B(c)=V]\ge1-\frac VQ.
\end{equation}
In the random expanded-key model, $\mathcal B(c)$ is a uniform $m\times V$ matrix for
every fixed nonzero $c$.  Moreover, the coordinate-direction matrices
$\mathcal B(e_1),\ldots,\mathcal B(e_O)$ are independent, so
\begin{equation}\label{eq:all-coordinate-directions-bad}
 \Pr\bigl[\rank\mathcal B(e_j)<V\text{ for every }j\bigr]
 =\tau_{m,V}(Q)^O.
\end{equation}
\end{proposition}

\begin{proof}
Every maximal minor of $\mathcal B(c)$ is homogeneous of degree $V$, so
Schwartz--Zippel gives \eqref{eq:random-direction-success}.  Under
Definition~\ref{def:ideal}, each $B_i$ is an independent uniform linear map.
This gives the fixed-direction law and independence across coordinate
directions.  The exact rank calculations appear in
\cref{app:directional-fallbacks}.
\end{proof}

The verified attack later uses a fresh random linear recombination to select
$V$ of the $m$ equations.  The supplement also gives a deterministic public
list of only $V(m-V)+1$ recombinations, together with exact list sizes and
failure probabilities for the official parameters
(\cref{app:directional-fallbacks}).

\subsection{An alternative recovery from a projective oil point}

The supplement gives a separate route that finds a nonzero point $x$ on the
oil space, considered only up to a common scalar
(\cref{sec:projective}).  Let $L/K$ be the field over which this point is
found.  From $x$, form the public matrix
\[
C_x=(P_1x\ \cdots\ P_mx)\in L^{(V+O)\times m}.
\]

\begin{lemma}[The common derivative kernel]\label{lem:kernel}
If $x\in\PP(\cO_L)$, then $\cO_L\subseteq\ker C_x^{\trans}$.  If
$\rank C_x=V$, equality holds.
\end{lemma}

\begin{proof}
For $u\in\cO_L$, \eqref{eq:isotropy} gives $u^{\trans}P_ix=0$ for every $i$.
If $\rank C_x=V$, the kernel of its transpose has dimension
$(V+O)-V=O$.
\end{proof}

When the derivatives at $x$ have rank $V$, the lemma recovers the entire oil
space as a kernel.  Row reduction, recovery of a basis over $K$, and the
checks from Propositions~\ref{prop:verified-key}
and~\ref{prop:signing-search}
then produce a verified signing key.  The attack does not need the oil space to
be unique: it keeps every exactly verified candidate with a signing witness.
\Cref{app:uniqueness} separately bounds the probability of a second
$K$-rational oil space.

\section{Why one root with invertible Jacobian is enough}\label{sec:local-geometry}

The exact solver used below comes with algebraic conditions on the complete
solution set.  At first sight, these conditions are stronger than what the
attack can prove.  The attack knows only that the planted root has an
invertible Jacobian.  This section shows that the solver may restrict its
attention to the set where the Jacobian is invertible, so singular roots and
components elsewhere are irrelevant.

Fix $c\in K^O$ and combine the $m$ public equations into $V$ equations using a
public matrix $R\in K^{V\times m}$.  Denote the resulting equations by
\[
 f=(f_1,\ldots,f_V)^{\mathsf T}:=R(q_{1,c},\ldots,q_{m,c})^{\mathsf T}
\]
and let $h$ be the determinant of their Jacobian:
\begin{equation}\label{eq:local-jacobian-polynomial}
 h(x):=\det J_f(x).
\end{equation}
Working only where $h\ne0$ is called \emph{localization at $h$}: algebraically,
we allow division by powers of $h$.  This removes every point where the
Jacobian is singular.  On the remaining set, the equations satisfy the two
technical conditions required by the solver.

\begin{lemma}[Solver conditions where the Jacobian is invertible]\label{lem:jacobian-localization}
Let $k$ be a perfect field, let $f_1,\ldots,f_V\in k[x_1,\ldots,x_V]$, and let $h=\det J_f$.  In the localization
\[
 S_h:=k[x_1,\ldots,x_V]_h,
\]
the sequence $f_1,\ldots,f_V$ is regular.  For every $1\le s\le V$, the ideal $(f_1,\ldots,f_s)S_h$ is radical.
\end{lemma}

Here regular means that the equations cut the dimension by one at each step.
Radical means that their common zeros carry no repeated algebraic
multiplicity.  The lemma applies these statements only on the set $h\ne0$.

The proof appears at the start of the finite-field solver appendix in the
technical supplement.  It applies the Jacobian criterion to each prefix and
then uses the complete-intersection criterion in $S_h$.

If $\mathcal B(c)$ has full column rank and $R\mathcal B(c)$ is invertible,
Proposition~\ref{prop:graph-direction-identity} gives
\begin{equation}\label{eq:h-at-planted}
 h(Gc)=(-2)^V\det(R\mathcal B(c))\ne0.
\end{equation}

\begin{corollary}[The solver conditions hold near the planted root]\label{cor:local-hypotheses}
If $\mathcal B(c)$ has full column rank and $R\mathcal B(c)$ is invertible,
then $Gc$ belongs to the set defined by the equations and inequality
\[
 \{f_1=\cdots=f_V=0,\ h\ne0\}.
\]
The input tuple $(f_1,\ldots,f_V,h)$ satisfies hypotheses (H1)--(H2) of the
solver of Gim\'enez et al.~\cite{Gimenez2026}, which accepts equations together
with the inequality $h\ne0$.  These conclusions are independent of every root
and algebraic component contained in $h=0$.
\end{corollary}

\begin{proof}
Equation \eqref{eq:h-at-planted} puts the planted root in the open set.
Lemma~\ref{lem:jacobian-localization} gives exactly the two algebraic
conditions required after localization.
\end{proof}

Thus the attack needs an invertible Jacobian at the planted root.  It does not
need the Jacobian to be invertible at every root or on every component of the
full system.


\section{Finding roots over the QR-UOV base field}\label{sec:finite-char-solver}

The equivalent-key attack needs one isolated root at which the Jacobian is
invertible.  It does not need all roots or a description of every algebraic
component.  We use the exact symbolic solver of Safey El Din and Schost for
this task~\cite{SafeySchost2018}.  The solver connects an easy system whose
roots are known to the target equations and follows each root as the
connecting parameter changes.  We call the path followed by one root a
\emph{branch}.  This section adapts the solver to characteristic $127$ and
states exactly what it returns.

\subsection{The roots sought by the solver}
All finite fields are perfect, which allows the Jacobian to detect repeated
roots.  Let $k$ be any perfect field and let
\[
 f=(f_1,\ldots,f_V)\in k[X_1,\ldots,X_V]^V
\]
be a square system of affine quadrics.  Write
\begin{equation}\label{eq:ns-locus}
 Z_{\rm ns}(f)
 :=\{x\in\overline k^V:f(x)=0,\ \det J_f(x)\ne0\}.
\end{equation}
Every point in $Z_{\rm ns}(f)$ is isolated.  B\'ezout's theorem gives
\begin{equation}\label{eq:C-C0}
 C:=2^V,
 \qquad
 C_0:=2^V+V2^{V-1}=2^{V-1}(V+2).
\end{equation}
Here $C$ bounds the number of isolated roots.  The solver represents the paths
from the easy system to the target as an algebraic curve in one connecting
parameter, and $C_0$ bounds the degree of that curve.  Recovering a rational
function of degree $C_0$ requires power-series coefficients through precision
$2C_0$.  This precision determines the main cost below.  The supplement gives
the multihomogeneous degree calculation that yields the formula for $C_0$.

The published algorithm starts from a product of univariate equations with
known simple roots and follows at most $C$ branches to the target equations.
It then discards outputs whose target Jacobian is not invertible.  The algorithm
also uses the separator introduced in \cref{sec:direct} to assign a scalar
value to each root.
The published proof uses a lower bound on the field characteristic only to
construct the starting equations and this separator.  We replace both
constructions.

\subsection{Why the solver works over the QR-UOV base field}

Safey El Din and Schost's published characteristic bound enters through two
finite families: the easy starting system and the separator
~\cite{SafeySchost2018}.  We replace the start by a Vandermonde product system
with exactly $2^V$ explicit simple roots, which requires only $2V$ distinct
field elements.  We replace the integer-indexed separator family by a uniform
linear form.  At most $C^2$ choices cause two roots to receive the same value,
so a uniform choice works with probability at least $1-C^2/|k'|$.

All remaining steps use exact field operations and power-series lifting.  They
do not divide by factorials or by any value that could vanish specifically in
characteristic $127$.  The technical supplement gives both constructions,
the lifting identities, and a step-by-step comparison with the published
algorithm.

\begin{theorem}[Roots with invertible Jacobian in characteristic $127$]\label{thm:finite-char-nonsingular}
Let $k$ be a finite field and let $f=(f_1,\ldots,f_V)$ be affine quadrics over
$k$, represented by an arithmetic circuit of size $L$.  Let $k'/k$ be a
finite extension containing at least $2V$ elements.  There is a randomized
algorithm with the following properties.
Here $\widetilde O$ suppresses factors that are polynomial in logarithms of
the displayed quantities.

\begin{enumerate}
\item Every returned point belongs to $Z_{\rm ns}(f)$.
\item Every point of $Z_{\rm ns}(f)$ is included in the output representation
with probability at least
\begin{equation}\label{eq:finite-char-success}
 1-\frac{C^2}{|k'|}.
\end{equation}
\item The arithmetic cost is
\begin{equation}\label{eq:published-slp-specialization}
 \softO\!\left(
 C C_0\bigl(L+2V+V^2\bigr)V
 \right)
\end{equation}
operations in $k'$.
\end{enumerate}
In particular, $|k'|\ge8C^2$ gives an inclusion probability of at least $7/8$
in one invocation.
\end{theorem}

\begin{proof}
Run the symbolic solver with the Vandermonde start and uniform
separator described above.  The construction and proof crosswalk in the
technical supplement show that the remaining steps introduce no further
characteristic bound.  The separator succeeds with probability at least
\eqref{eq:finite-char-success}; on that event the published specialization
recovers every isolated target root with invertible Jacobian.  Its final \textsc{Clean}
step removes every point with singular target Jacobian, even after a bad
separator.

For $V$ quadrics in one block, the degree quantities in the published
Proposition~5 are exactly $C$ and $C_0$ from \eqref{eq:C-C0}, and the sum of
degrees is $2V$.  The new starting system costs $O(V^2C)$ and does not change
the dominant lifting term.  Applying that proposition gives
\eqref{eq:published-slp-specialization}.
\end{proof}

\begin{remark}[The attack needs validity, not a complete root list]\label{rem:one-sided-enough}
The original solver paper notes that it cannot cheaply certify that a generic
output contains every root.  QR-UOV does not need such a certificate.  We
check every original fixed-direction equation, require the root coordinates to lie
in $K$, complete the graph, and verify the unchanged public map.  An omitted
planted root causes a restart.  A candidate that is not a valid root or key is
rejected.
\end{remark}

\subsection{Complete equivalent-key attack}

We now combine the fixed-direction equations, the root solver, graph completion,
and the public signing test into one attack.

For a random $V\times m$ matrix $R$ over $K$, put
\[
 f_{c,R}:=R(q_{1,c},\ldots,q_{m,c})^{\mathsf T}.
\]
If $\mathcal B(c)$ has rank $V$, then $R\mathcal B(c)$ is a uniform $V\times V$ matrix.  Write
\begin{equation}\label{eq:mu-square}
 \mu_V(Q):=\prod_{j=1}^V(1-Q^{-j}),
 \qquad Q=|K|=127^3.
\end{equation}

Choose $K'=\F_{Q^r}$ with
\begin{equation}\label{eq:r-choice}
 r:=\min\{s\ge1:Q^s\ge8\cdot2^{2V}\}.
\end{equation}
The values are
\begin{equation}\label{eq:r-values}
 r=6,\qquad8,\qquad10
\end{equation}
at Levels I, III, and V.

\begin{theorem}[Verified equivalent-key recovery]\label{thm:verified-directional-recovery}
Let a fixed expanded \QRUOV{} public key have a valid graph decomposition $G$.  Suppose
\begin{enumerate}
\item some maximal minor of $\mathcal B(c)$ is a nonzero polynomial in $c$; and
\item the recovered planted decomposition has a nonzero signing determinant polynomial $\Delta$ from Proposition~\ref{prop:signing-search}.
\end{enumerate}
Then the following algorithm eventually returns another signing key for the
same public key.  It never returns an incorrect key and repeats with fresh
randomness when an attempt fails.

\begin{enumerate}
\item Sample $c\leftarrow K^O$ and $R\leftarrow K^{V\times m}$.
\item Run \cref{thm:finite-char-nonsingular} on $f_{c,R}$ over $K'$ from \eqref{eq:r-choice}.
\item Keep only roots that solve all $m$ original fixed-direction equations and
whose coordinates lie in $K$.
\item For every retained root $x$, run the graph completion of
Proposition~\ref{prop:graph-direction-completion}.  Require all graph
identities, the inverse pull-back, exact public-map factorization, and a
signing witness.
\item Return the first accepted key; otherwise restart with fresh public coins.
\end{enumerate}

One invocation includes the planted branch with probability at least
\begin{equation}\label{eq:directional-invocation-success}
 p_{\rm dir}
 \ge
 \left(1-\frac VQ\right)
 \mu_V(Q)
 \left(1-\frac{2^{2V}}{Q^r}\right).
\end{equation}
The expected numbers of full solver invocations are at most
\begin{equation}\label{eq:directional-restart-values}
 1.0000262,
 \qquad1.0000561,
 \qquad1.0202201.
\end{equation}
Every returned key is correct, regardless of whether failed solver invocations
included every root.
\end{theorem}

\begin{proof}
By Proposition~\ref{prop:graph-direction-density}, a random direction gives a
full-rank derivative at the planted root with probability at least $1-V/Q$.
Conditional on that event, $R\mathcal B(c)$ is a uniform square matrix, so it
is invertible with probability $\mu_V(Q)$.  Proposition~\ref{prop:graph-direction-identity}
then makes the planted point an isolated root of $f_{c,R}$ with invertible
Jacobian.  Theorem~\ref{thm:finite-char-nonsingular} includes it with
probability at least $1-2^{2V}/Q^r$.

The filters for the original equations and for coordinates in $K$ retain the
planted point.  Proposition~\ref{prop:graph-direction-completion} reconstructs
$G$, and Propositions~\ref{prop:verified-key} and~\ref{prop:signing-search}
supply the exact factorization and signing certificate.  Conversely, the
algorithm rejects a candidate that fails any original equation, graph
completion equation, graph identity, pull-back identity, factorization check,
or oil-matrix test.  Thus every accepted key is correct.  Retrying with fresh
randomness and using \eqref{eq:directional-invocation-success} gives the
numerical expectations above.
\end{proof}

\begin{corollary}[Probability that all public choices fail]\label{cor:directional-public-random}
In the random expanded-key model, the probability that every coordinate
direction has deficient planted derivative is
$\tau_{m,V}(Q)^O$, with base-two logarithms
\[
 -1132.166907,\qquad-1635.352198,\qquad-2935.247544.
\]
The probability that the deterministic signing search over pairs of basis
matrices fails has base-two logarithm below
\[
 -544.820,\qquad-796.276,\qquad-1068.687.
\]
Outside their union, a deterministic scan of coordinate directions and
signing witnesses includes a successful choice.  Only the internal root solver
requires randomness.
\end{corollary}

\begin{proof}
Combine Propositions~\ref{prop:graph-direction-density}
and~\ref{prop:signing-search}.  The two exceptional events need not be
independent, so a union bound suffices.
\end{proof}

\section{Concrete cost of the symbolic solver}\label{sec:concrete-solver}

The preceding theorem gives an asymptotic cost, but a concrete estimate also
depends on how the solver represents and multiplies large polynomials.  This
section states the resulting estimate first.  It then explains the choices
that change the final number: evaluating all quadratic equations together,
using extension fields with sparse defining polynomials, distinguishing only
the root that may yield a forgery, and testing exactly whether that root lies
in the base field.  The section ends with the success probability and the
complete Level-I cost.

\subsection{Resulting conditional cost estimate}\label{sec:local-separator-upgrade}

Let $q=127$, $a=m-4$, $C=2^a$, and $C_0=2^{a-1}(a+2)$ as in the direct
attack.  We compute the cost for several extension fields whose sparse
defining polynomials are proved irreducible, then select the least expensive
degree.

\begin{theorem}[Conditional Boolean-gate estimate]\label{thm:headline}
Assume the random expanded-key model, the degree bound for the curve followed
by the solver, the polynomial multiplication costs stated below, and the
specification's Boolean-gate convention.  Require the separator to distinguish
only a root that passes the forgery tests, as justified in \cref{sec:local},
and use the exact base-field test in \cref{sec:frob}.  After accounting for
failed attempts, the estimated costs are those in \cref{tab:headline}.  In
particular, Level I costs
\[
 2^{140.383995063}
\]
Boolean gates, leaving $2.616004937$ bits below the category target.
\end{theorem}

\begin{table}[ht]
\centering
\caption{Estimated base-two logarithms of Boolean-gate costs.  ``All roots''
requires the separator to distinguish every root.  ``Accepted root'' requires
it only for a root that passes the forgery tests.}
\label{tab:headline}
\begin{tabular}{@{}lrrrrr@{}}
\toprule
Level & Original & All roots & Accepted root & Target & Margin\\
\midrule
I   & 151.639 & 141.126 & 140.384 & 143 & $+2.616$\\
III & 203.507 & 191.630 & 191.027 & 207 & $+15.973$\\
V   & 260.389 & 247.637 & 247.125 & 272 & $+24.875$\\
\bottomrule
\end{tabular}
\end{table}

The least expensive Level-I row uses $K'=\F_{127^8}$.  The Rabin
irreducibility test certifies the following defining polynomial:
\[
 X^8+X^4-1.
\]
The nearby degree-$9$ choice $X^9+X^2+1$ gives $140.418965$, which provides a
cross-check on the optimization.

\subsection{Batching equation evaluation and polynomial multiplication}\label{sec:flat-lifting}

The $V$ dense quadrics share one vector of quadratic monomials.  Evaluating
that vector once and using it for all $V$ equations removes a redundant factor
from the arithmetic-circuit cost.  To follow the roots, the solver stores
truncated power series in a parameter $T$.  A second variable $U$ labels the
roots through the separator.  At precision $k$, these coefficients lie in
\[
 \mathcal A_k=K'[T,U]/(T^k,Q_k(U)),
 \qquad \deg_U Q_k\le C.
\]
Here $Q_k$ is the monic polynomial whose roots encode the current root values.
We encode the pair of exponents $(i,j)$ as one exponent, using radix $2C-1$.
This converts multiplication in $\mathcal A_k$ into three ordinary polynomial
products, including reduction modulo $Q_k$.  Let $\Mmul(t)$ denote the cost of
multiplying polynomials of length $t$.  Across all stages that double the
precision, one product in $\mathcal A_k$ is bounded by
\begin{equation}\label{eq:flat-schedule-bound}
 3\Mmul(8CC_0).
\end{equation}
The factor $8$ accounts for the final precision $2C_0$.  The technical
supplement gives the batched circuit, the exponent encoding, reduction modulo
$Q_k$, and the proof of the resulting asymptotic cost
$\widetilde O(V^{\omega+1}4^V)$.

\subsection{Sparse working fields}\label{sec:sparse-fields}

The solver performs arithmetic in an extension of $\F_{127}$.  A sparse
defining polynomial reduces the number of base-field operations needed after
each multiplication.  The direct solver may use any extension degree.  The
equivalent-key solver starts over $K=\F_{127^3}$, so its degree over
$\F_{127}$ must be divisible by three.  Rabin tests certify sparse irreducible
polynomials of degrees $15,22,30$ for direct forgery and $18,24,30$ for
equivalent-key recovery.  The refined direct estimate separately selects
degrees $8,12,16$.

For an extension degree $d$, let $(M_d,A_d)$ be the audited counts of base-field
multiplications and additions in the selected recursive Karatsuba circuit, and
let $s_d$ count the nonleading terms of the sparse modulus.  One multiplication
in the working field is charged by the following upper bound:
\begin{equation}\label{eq:sparse-extension-envelope}
 \mathsf G^{\rm sparse}_{127^d}
 =M_d\Gmul(\log_2 127)
  +\bigl(A_d+s_d(d-1)\bigr)4\log_2 127.
\end{equation}
Here $\Gmul(b)=2b^2+b$ is the specification's gate charge for multiplying two
$b$-bit values.  The supplement gives every defining polynomial, its
irreducibility certificate, the recursive operation count, and a direct
multiplication cross-check.

\subsection{Collisions among rejected roots do not raise the degree}
\label{sec:charpoly}

This subsection explains why only the root that eventually passes the forgery
tests needs a unique scalar label.  Collisions among roots that will be
rejected need not increase the solver's cost.

As the connecting parameter $t$ changes, the solver represents all moving
roots at once in a finite-dimensional algebra.  After removing components that never reach the
target equations, the corresponding curve has degree at most $D=C_0$.
Introduce formal coefficients $\Lambda_j$ and assign each root the value
$\sum_j\Lambda_jx_j$.  The characteristic polynomial
$\chi_\Lambda(t,U)$ has these values as its roots.  A concrete linear form
$\lambda$ may assign the same value to two roots, but the following lemma
shows that such a collision does not increase the degree in $t$.

\begin{lemma}[A collision does not increase the degree]
\label{lem:charpoly}
For every $\lambda\in K'^a$, including a nonseparating choice,
\[
 q_\lambda(t,U)=\chi_\Lambda(t,U)|_{\Lambda=\lambda}
\]
has a common-denominator representation whose numerator and denominator have
$t$-degree at most $D$.
\end{lemma}

\begin{lemma}[Recovering coordinates by differentiation]\label{lem:polar}
The interpolation numerator for coordinate $x_j$ is
\[
 v_j(t,U)=-
 \left.\frac{\partial\chi_\Lambda(t,U)}{\partial\Lambda_j}
 \right|_{\Lambda=\lambda}.
\]
It has the same degree bound.  The numerator for any fixed linear combination
of the coordinates also has degree at most $D$.  The attack additionally uses
the rational function $H_0/(\delta+\lambda)$ to test the final target equation.
After adjoining this function as another coordinate, its numerator has degree
at most $2D$.
\end{lemma}

Schost's parametric degree theorem gives these statements first for the generic
linear form~\cite{Schost2003}.  Specializing $\Lambda$ can create repeated
factors, but it cannot increase the degree in $t$ because the common
denominator is independent of $\Lambda$.  The attack may therefore keep the
full characteristic polynomial and require $\lambda$ to distinguish only the
root it accepts.  The additional rational test gives reconstruction precision
$4C_0+1$.  The concrete cost still assumes that Schost's degree theorem
applies after removing the irrelevant components of the exact attack curve
and after adding this rational test.

\subsection{Only the accepted root needs a unique label}\label{sec:local}

When the connecting parameter reaches the target equations, the interpolated
polynomials may share a power of that parameter.  Cancel this common power.
Let $\bar q$ be the resulting univariate polynomial whose roots label target
candidates, and let $\bar v_b$ be the corresponding numerator for a scalar
value $b$.  The exponent of the lowest nonzero power of the connecting
parameter is called its \emph{valuation}.  Some branches may diverge as the
target is approached.  The linear form $\lambda$ must preserve the lowest
nonzero power on those branches
so that canceling the common power does not remove a finite target root.

\begin{lemma}[Recovering one uniquely labeled root]\label{lem:simple}
If $\lambda$ preserves the minimum valuation of every branch escaping to
infinity and $\tau$ is a simple root of $\bar q$, then exactly one finite
represented point $P$ satisfies $\lambda(P)=\tau$, and
\[
 b(P)=\frac{\bar v_b(\tau)}{\bar q'(\tau)}.
\]
No separation condition is needed at any other root value.
\end{lemma}

For a fixed usable target point, at most $C-1$ linear conditions on a uniform
$\lambda\in(K')^a$ can violate the valuation condition or collide with that
point.  Hence
\begin{equation}\label{eq:lambdafail}
 \Pr[\text{bad target separator}]\le\frac{C-1}{|K'|}.
\end{equation}
The rational test $H_0/(\delta+\lambda)$ uses a random scalar $\delta$ in its denominator.
After rejecting shifts that vanish at a known starting root, its remaining
denominator failure probability is
\begin{equation}\label{eq:cfail}
 \Pr[\text{bad target denominator}]\le\frac{C}{|K'|-C}.
\end{equation}
The earlier analysis required one linear form to distinguish every pair among
at most $C$ roots, which cost $C^2/|K'|$.  The supplement proves the weaker
condition above and both failure bounds.

\subsection{One power test identifies base-field roots}\label{sec:frob}

The solver works over $K'=\F_{q^d}$, but a QR-UOV forgery needs coordinates
in the base field $\F_q$.  For a candidate point $P$, set
$b_1(P)=\lambda^{q^{d-1}}\cdot P$.

\begin{lemma}[Exact base-field test for a simple root]\label{lem:exactfrob}
If $\tau=\lambda(P)$ is a simple root of $\bar q$, then
\[
 b_1(P)^q=\tau\quad\Longleftrightarrow\quad P\in\F_q^a.
\]
\end{lemma}

Indeed, the equality says $\lambda(P^q)=\lambda(P)$.  Raising every coordinate
to its $q$th power maps target roots to target roots.  If $P^q\ne P$, these two
distinct roots would receive the same value under $\lambda$, making $\tau$ a
repeated root.  The test therefore accepts exactly the candidates in
$\F_q^a$ and needs no second linear combination of the coordinates.

\subsection{Root success probability}\label{sec:rootbound}

Let $M$ count all off-oil affine roots and $N$ those with invertible Jacobian.
The exact first and second moments, together with the pairing $z,-z$, give
\[
 \Pr[\text{usable projective root with full-rank Jacobian}]
 \ge \frac{\mu\pi}{2}-\frac{\mu^2}{8}
   +\frac{(q-1)\mu q^{-r}}8,
\]
where $\mu\ge1-q^{-1}$ and
$\pi=\prod_{j=1}^{r-1}(1-q^{-j})$.  Substitution yields
\[
 p_{\rm root}=0.3690869822\ldots.
\]
Using instead the earlier proved bound
$0.326336998767509\ldots$ raises the Level-I total only to $140.531461$ bits,
still $2.469$ bits below target.  The supplement retains the moment
calculation.

\subsection{Level-I cost calculation}

At the selected extension degree $d=8$, one attempt succeeds with probability
at least $0.353997$.  The cost therefore includes an expected $2.824886$
attempts.  The failure probabilities for distinguishing the accepted root and
for the denominator in its rational test are at most $0.016637$ and
$0.016918$.  We charge the full power-series lifting schedule, two sets of
scalar interpolation data, one coordinate reconstruction, factorization, and
public verification.  We also add the cost of $64$ long polynomial products
for lower-order work.  The total is
\[
 2^{140.383995063}
\]
Boolean gates.  The result stays below $2^{143}$ when the
cost assigned to rational reconstruction ranges from $8$ to $128$ times its
baseline.  It also stays below target if the assumed cost of the long
polynomial products is multiplied by eight, but not by sixteen.  The technical
supplement gives every component and both sensitivity tables.


\section{Evidence and limitations}\label{sec:evidence-limitations}

The two reductions and the fixed-output theorem are exact algebraic results.
The same is true of the identities used to recover one distinguished root and
to test whether its coordinates lie in the base field.  None of these
statements depends on the Boolean-gate cost model.  The attack checks every
candidate against the original equations or the unchanged public verifier.
A failed solver call may force another attempt, but it cannot make the
verifier accept an invalid signature or key.

We implemented the direct-forgery attack on smaller parameters.  The program
starts from expanded public matrices, a message, and a salt.  It produces a
forged signature and verifies it using only public data.  Six instances over
$\F_{13}$ and one over $\F_{29}$ completed successfully, and a separate public
replay reproduced each result.  The program keeps the QR-UOV quadratic map but
simplifies the conversion from seeds, the byte encoding, and the hash
interface.  It also uses straightforward power-series arithmetic rather than
the fast multiplication assumed by the full cost estimate.  The experiment
therefore checks the attack from the expanded key through verification.  It
does not test the official input encoding or the full-parameter running time
(\cref{app:reduced-implementation}).

The full-parameter cost estimate still depends on four unproved modeling
steps:
\begin{enumerate}
\item Schost's degree theorem must apply to the exact attack curve after
irrelevant components are removed and the rational target test is added.
\item The long polynomial products must admit a complete implementation with
the cost assumed in \cref{sec:flat-lifting}.
\item The random expanded-key model must accurately describe the official
seeded key expansion.
\item The implementation must move and store its polynomial coefficients, a
cost omitted by the Boolean-gate count and already prohibitive at the stated
memory sizes.
\end{enumerate}
The Level-I margin is only $2.616$ bits, so each assumption can affect the
headline conclusion.  The technical supplement gives the full validation
record and the commands used to reproduce the smaller experiments.



\begin{table}[H]
\centering
\small
\caption{Status of the main claims.}
\label{tab:claim-status-final}
\begin{tabularx}{\textwidth}{@{}p{0.31\textwidth}X@{}}
\toprule
Claim & Status \\
\midrule
Direct system in $m-4$ variables & Exact for every nonzero target once an
initial common zero is supplied.  Chevalley--Warning guarantees this zero in
every seven-dimensional starting subspace, and fixing the common scale removes
one variable. \\
Initial root lies outside the oil space & Probability bound for every fixed
key over the attack's public randomness.  This statement does not average over
key generation. \\
Construction of the attack subspace & Deterministic and finite after the initial root, using linear algebra and exhaustive search in $\PP^6(\F_q)$. \\
Probability of a usable direct root & Proved lower bound in the independent-graph random expanded-key model. \\
A root with full-rank derivative determines another signing key & Exact for every expanded key that passes the public rank, graph, and signing checks. \\
Returned signature or key is valid & Exact; every candidate is verified against the unchanged public map. \\
Concrete cost rows & Calculations for symbolic algorithms, not measurements from implemented circuits. \\
Full-parameter memory & Exceeds physical storage for the symbolic representations analyzed here. \\
Fixed SHAKE or AES expansion & Verification is exact for every key.  The
success probabilities transfer only under the stated random-function or
ideal-cipher models, with explicit terms for finite output length and seed
collisions. \\
\bottomrule
\end{tabularx}
\end{table}

\section{Cost, memory, and parameter consequences}\label{sec:security}

\subsection{Work unit and scope}
The QR-UOV specification expresses classical attack costs as Boolean-gate
exponents.  A reported cost $c$ therefore means approximately $2^c$ Boolean
gates.  The specification charges multiplication of two $b$-bit field
elements by
\begin{equation}\label{eq:spec-gmul-new}
 \Gmul(b)=2b^2+b
\end{equation}
Boolean gates~\cite{QRUOVSpec}.  We use the same unit.  We upper-bound every
extension-field operation by one multiplication using the explicit
divide-and-conquer Karatsuba circuit in \eqref{eq:sparse-extension-envelope}.

The main solver cost comes from multiplying the truncated polynomials that
represent the moving roots.  For a system with $a$ equations in $a$ variables,
\[
 C=2^a,
 \qquad C_0=2^{a-1}(a+2),
\]
the solver reconstructs through precision $2C_0$.  For the concrete numbers,
we assume that multiplying polynomials of length $t$ costs
\[
 \Mmul(t)=t\log_2t\log_2\log_2t.
\]
The complete power-series calculation is then charged by
\begin{equation}\label{eq:unified-flat-cost}
 3(a^3+a^2)\Mmul(8CC_0)
\end{equation}
operations in the chosen working field.  The direct-forgery row uses
$a=m-4$.  Its cost includes the probabilities that the initial subspace avoids
the oil space, that a usable root exists, that the chosen chart contains it,
and that the separator identifies it.  The equivalent-key row uses $a=V$ and
includes the corresponding full-rank and separator probabilities.  The
calculation also includes the lower-order costs of constructing the equations,
filtering roots, factoring polynomials, completing the graph, restoring the
scale, and verifying the final output.  These terms do not change the displayed
decimals.

These are calculations for symbolic algorithms, not measurements from an
implemented circuit.  In particular, the formula for $\Mmul$ assigns a unit
constant to an asymptotic multiplication bound.  The algebraic reductions,
field-multiplication circuits, chosen fields, and success probabilities do not
depend on this convention.  A practical running time does.

\subsection{The two attacks}
The following table reports the estimated cost exponent and its distance from
the category target for both attacks.
\begin{table}[t]
\centering
\small
\setlength{\tabcolsep}{5pt}
\caption{Estimated base-two logarithms of classical Boolean-gate costs.  A
positive margin means that the estimated attack cost is below the category
target.}
\label{tab:combined-costs}
\begin{tabular}{@{}lrrrrrr@{}}
\toprule
& \multicolumn{2}{c}{Level I} & \multicolumn{2}{c}{Level III} & \multicolumn{2}{c}{Level V}\\
Attack & Cost & Margin & Cost & Margin & Cost & Margin\\
\midrule
Direct forgery, original solver
 &151.639&$-8.639$&203.507&$3.493$&260.389&$11.611$\\
Direct forgery, accepted-root solver
 &140.384&$2.616$&191.027&$15.973$&247.125&$24.875$\\
Equivalent-key recovery
 &154.826&$-11.826$&206.177&$0.823$&260.857&$11.143$\\
\midrule
Best attack
 &\textbf{140.384}&$\mathbf{2.616}$
 &\textbf{191.027}&$\mathbf{15.973}$
 &\textbf{247.125}&$\mathbf{24.875}$\\
\midrule
Category target &143&&207&&272&\\
\bottomrule
\end{tabular}
\end{table}

Both direct rows solve systems in $50,74,101$ variables.  The second row is
lower because the separator needs to distinguish only an accepted root, one
$q$th-power equality identifies base-field roots, and coordinates are reconstructed
only after a candidate passes the scalar tests.  The selected extension
degrees are $8$, $12$, and $16$.  The equivalent-key row uses the independent
attack from \cref{sec:recovery}.  The smaller implementation executes the
refined direct attack, but the full-parameter polynomial calculation still
dominates both time and memory.

For context, the QR-UOV team's January 2026 response reports classical
Hashimoto estimates of $158,217,285$ bits for the same three public
dimensions~\cite{QRUOVResponse2026}.  The new direct rows are lower by
approximately $6.36,13.49,24.61$ bits.  The estimates use different models.
The team assumes the usual semi-regular behavior, a standard heuristic for
the growth of Gr\"obner-basis computations on generic-looking equations.
Our row instead uses the exact symbolic solver from
\cref{sec:finite-char-solver} and its proved success probability.

\subsection{What follows at each level}
The refined direct rows lie below the three targets by $2.616$, $15.973$, and
$24.875$ bits.  At Level I this is a candidate break only under the stated
model.  The polynomial representation cannot fit in physical memory, the fast
multiplication schedule has not been implemented at full size, and the degree
theorem has not been checked for every detail of the exact attack curve.  The
larger Level-III and Level-V margins, together with the independent
equivalent-key estimates, support changing those parameter sets.  The smaller
implementation checks that the direct-attack steps work together.  It does not
make any full-parameter row practical.

\subsection{Why changing only the vinegar count is not a repair}
The official parameter sets sit on the direct side of the crossing between the
two reductions:
\begin{center}
\begin{tabular}{@{}lrrrr@{}}
\toprule
Level & $V$ & $m-4$ & Key-recovery variables & Direct variables\\
\midrule
I&52&50&52&50\\
III&76&74&76&74\\
V&102&101&102&101\\
\bottomrule
\end{tabular}
\end{center}
Increasing $V$ raises the key-recovery dimension but leaves the direct-forgery
dimension unchanged.  Decreasing $V$ eventually makes the
equivalent-key system no larger than $m-4$.  Thus, with $m$ fixed, changing
only $V$ always leaves at least one of the two systems at or below this size.

The next table applies the same conservative cost calculation to nearby
parameters.  It is not a complete replacement proposal.  Any proposal must
also recompute other attacks, failure probabilities, key sizes, and
implementation costs.  The table isolates the consequence proved here:
raising only $V$ leaves direct forgery unchanged, while raising only $m$
leaves the equivalent-key system unchanged.

\begin{table}[t]
\centering
\small
\caption{Smallest $(V,m)$ pairs with $m$ divisible by three in the cost calculation
that put both displayed attacks at the category target, or $16$ bits above the
target.  This is not a complete parameter proposal.}
\label{tab:combined-retuning}
\begin{tabular}{@{}lrrrrrr@{}}
\toprule
& \multicolumn{2}{c}{Current} & \multicolumn{2}{c}{Reach target} & \multicolumn{2}{c}{Reach target $+16$}\\
Level & $V$ & $m$ & $V$ & $m$ & $V$ & $m$\\
\midrule
I&52&54&52&54&57&60\\
III&76&78&78&81&87&90\\
V&102&105&108&111&117&120\\
\bottomrule
\end{tabular}
\end{table}

The table enforces $m\equiv0\pmod3$ and the condition $V\ge m-3$ needed by the
direct reduction.  Level I already reaches its category target in this model,
but placing the attack cost $16$ bits above the target requires increasing
both dimensions.  At Levels III and V, even reaching the target requires
increasing both $V$ and $m$.  A design change that invalidates one reduction
may cost less than increasing both dimensions.

\subsection{The symbolic representation does not fit in memory}\label{subsec:memory-new}
The desired outputs are small.  Direct forgery needs one vector in $\F_q^n$,
and equivalent-key recovery needs a $V\times O$ graph.  The solver, however,
stores a large array of polynomial coefficients while following all candidate
roots to the target equations.  One such array at the final precision requires
approximately
\[
 2^{49.41}\ \mathrm{EiB},
 \qquad2^{98.51}\ \mathrm{EiB},
 \qquad2^{153.40}\ \mathrm{EiB}
\]
for the direct systems.  The corresponding equivalent-key arrays require
about $2^{53.73}$, $2^{102.68}$, and $2^{155.41}$ EiB.  These quantities exceed
physical storage by many orders of magnitude.

No full-parameter forgery or equivalent-key recovery is reported.  The paper
therefore treats the algebraic reductions, modeled running time, storage, and
verification as separate claims.  The exact reductions remain relevant to
parameter selection even though this symbolic representation is not a
practical attack implementation.

\section{Conclusion}\label{sec:conclusion}
Changing only the number of vinegar variables cannot make both attacks harder.
For every vinegar count and every nonzero verification target, either direct
forgery or equivalent-key recovery produces at most $m-4$ nonlinear equations
in the same number of variables.  At the recommended parameter sets, the
direct systems have $50$, $74$, and $101$ variables.  When the vinegar count
is smaller, the equivalent-key route supplies the smaller system.  This
algebraic statement is exact and does not depend on
the cost of a particular equation solver.

We also estimate the work needed to solve these systems.  The adapted symbolic
solver needs its linear form to distinguish only a root that passes the
forgery tests, and one exact field-membership test determines whether that
root lies in the base field.  Under the stated Boolean-gate and polynomial-arithmetic model, the
estimated costs lie below the three category targets.  The Level-I margin is
$2.616$ bits.

The smaller implementation produces forged signatures from expanded public
keys and verifies them with public data.  It confirms the algebraic steps, but
not the full-parameter cost.  The concrete estimates still rely on the stated
degree theorem, polynomial multiplication cost, and random expanded-key model.
They also require more memory than any physical system can provide.  The firm
design conclusion is therefore narrower: a parameter revision must increase
the output dimension as well as the vinegar dimension, or change the
construction so that one of the two reductions no longer applies.

\begin{thebibliography}{99}
\fontsize{8.5pt}{9.6pt}\selectfont
\setlength{\itemsep}{0pt}
\setlength{\parsep}{0pt}

\bibitem{QRUOVOriginal}
H.~Furue, Y.~Ikematsu, Y.~Kiyomura, and T.~Takagi.
\newblock A new variant of Unbalanced Oil and Vinegar using quotient ring: QR-UOV.
\newblock Cryptology ePrint Archive, Paper 2020/1243, 2020.

\bibitem{LiGuoDing2025}
P.~Li, H.~Guo, and J.~Ding.
\newblock Unified approach to UOV-like multivariate signature schemes.
\newblock Cryptology ePrint Archive, Paper 2025/1778, 2025.

\bibitem{QRUOVSpec}
H.~Furue, Y.~Ikematsu, F.~Hoshino, T.~Takagi, H.~Kosuge, K.~Yamakoshi, R.~Akiyama, S.~Nakamura, S.~Orihara, and K.~Kinjo.
\newblock \emph{QR-UOV Specification Document (Round 2 submission)}.
\newblock NIST Additional Digital Signatures Project, 5 February 2025.

\bibitem{KPG1999}
A.~Kipnis, J.~Patarin, and L.~Goubin.
\newblock Unbalanced Oil and Vinegar signature schemes.
\newblock In \emph{EUROCRYPT 1999}, pp.~206--222, 1999.

\bibitem{KipnisShamir1998}
A.~Kipnis and A.~Shamir.
\newblock Cryptanalysis of the Oil and Vinegar signature scheme.
\newblock In \emph{CRYPTO 1998}, pp.~257--266, 1998.

\bibitem{BraekenWolfPreneel2005}
A.~Braeken, C.~Wolf, and B.~Preneel.
\newblock A study of the security of Unbalanced Oil and Vinegar signature schemes.
\newblock Cryptology ePrint Archive, Paper 2004/222, 2004.

\bibitem{FieldLifting}
W.~Beullens and B.~Preneel.
\newblock Field lifting for smaller UOV public keys.
\newblock In \emph{INDOCRYPT 2017}, pp.~227--246, 2017.

\bibitem{BACUOV}
A.~Szepieniec and B.~Preneel.
\newblock Block-anti-circulant Unbalanced Oil and Vinegar.
\newblock In \emph{SAC 2019}, pp.~574--588, 2019.

\bibitem{BACUOVAttack}
H.~Furue, K.~Kinjo, Y.~Ikematsu, Y.~Wang, and T.~Takagi.
\newblock A structural attack on block-anti-circulant UOV at SAC 2019.
\newblock In \emph{PQCrypto 2020}, pp.~323--339, 2020.

\bibitem{DingEtAl2008}
J.~Ding, B.-Y.~Yang, C.-O.~Chen, M.-S.~Chen, and C.-M.~Cheng.
\newblock New differential-algebraic attacks and reparametrization of Rainbow.
\newblock In \emph{ACNS 2008}, pp.~242--257, 2008.

\bibitem{NISTRound3}
National Institute of Standards and Technology.
\newblock Nine candidates advance to the third round of the Additional Digital Signatures for the PQC standardization process.
\newblock NIST CSRC, 14 May 2026.

\bibitem{NISTIR8610}
G.~Alagic et al.
\newblock Status report on the second round of the Additional Digital Signature Schemes for the NIST Post-Quantum Cryptography Standardization Process.
\newblock NIST IR~8610, May 2026.

\bibitem{Benoist}
O.~Benoist.
\newblock Degr\'es d'homog\'en\'eit\'e de l'ensemble des intersections compl\`etes singuli\`eres.
\newblock \emph{Annales de l'Institut Fourier}, 62(3):1189--1214, 2012.

\bibitem{vdHL}
J.~van der Hoeven and G.~Lecerf.
\newblock On the complexity exponent of polynomial system solving.
\newblock \emph{Foundations of Computational Mathematics}, 21(1):1--57, 2021.

\bibitem{GiustiLecerfSalvy2001}
M.~Giusti, G.~Lecerf, and B.~Salvy.
\newblock A Gr\"obner free alternative for polynomial system solving.
\newblock \emph{Journal of Complexity}, 17(1):154--211, 2001.
\newblock doi:10.1006/jcom.2000.0571.

\bibitem{Gimenez2026}
N.~Gim\'enez, J.~Heintz, G.~Matera, L.~M.~Pardo, M.~P\'erez, and M.~Privitelli.
\newblock A Kronecker algorithm for locally closed sets over a perfect field.
\newblock arXiv:2512.14888v2, 5 June 2026.

\bibitem{Beullens}
W.~Beullens.
\newblock Improved cryptanalysis of UOV and Rainbow.
\newblock In \emph{EUROCRYPT 2021, Part I}, pp.~348--373, 2021.

\bibitem{FurueIkematsu}
H.~Furue and Y.~Ikematsu.
\newblock A new security analysis against MAYO and QR-UOV using rectangular MinRank attack.
\newblock In \emph{IWSEC 2023}, pp.~101--116, 2023.

\bibitem{Suzuki}
T.~Suzuki, H.~Furue, T.~Ito, S.~Nakamura, and S.~Uchiyama.
\newblock An extended rectangular MinRank attack against UOV and its variants.
\newblock In \emph{IWSEC 2025}, LNCS~16208, pp.~111--130, Springer, 2026.

\bibitem{FurueIkematsu2026}
H.~Furue and Y.~Ikematsu.
\newblock Key recovery attacks on UOV using $p^\ell$-truncated polynomial rings.
\newblock Cryptology ePrint Archive, Paper 2026/298, 2026.

\bibitem{JustGuess}
A.~May, M.~Ostuzzi, and H.~Ressler.
\newblock Just Guess: Improved (quantum) algorithm for the underdetermined MQ problem.
\newblock In \emph{EUROCRYPT 2026, Part IV}, LNCS~16544, pp.~334--362, 2026.

\bibitem{Pebereau}
P.~P\'ebereau.
\newblock One vector to rule them all: Key recovery from one vector in UOV schemes.
\newblock In \emph{PQCrypto 2024}, pp.~92--108, 2024; ePrint 2023/1131.


\bibitem{BaurStrassen}
W.~Baur and V.~Strassen.
\newblock The complexity of partial derivatives.
\newblock \emph{Theoretical Computer Science}, 22(3):317--330, 1983.
\newblock doi:10.1016/0304-3975(83)90110-X.

\bibitem{DAndreaDickenstein}
C.~D'Andrea and A.~Dickenstein.
\newblock Explicit formulas for the multivariate resultant.
\newblock \emph{Journal of Pure and Applied Algebra}, 164(1--2):59--86, 2001.
\newblock doi:10.1016/S0022-4049(00)00145-6.

\bibitem{JeronimoEtAl}
G.~Jeronimo, T.~Krick, J.~Sabia, and M.~Sombra.
\newblock The computational complexity of the Chow form.
\newblock \emph{Foundations of Computational Mathematics}, 4(1):41--117, 2004.
\newblock doi:10.1007/s10208-002-0078-2.

\bibitem{NeigerEtAl}
V.~Neiger, B.~Salvy, E.~Schost, and G.~Villard.
\newblock Faster modular composition using two relation matrices.
\newblock In \emph{ISSAC 2026}, pp.~315--324, 2026; arXiv:2601.17422.

\bibitem{DegreeSmoothing}
S.~Jaques, L.~Ran, S.~Samardjiska, and M.~Seitner.
\newblock Smoothing the degree of regularity for polynomial systems.
\newblock Cryptology ePrint Archive, Paper 2026/408, 2026.

\bibitem{F4Tracer}
R.~Kouba, V.~Neiger, and M.~Safey El Din.
\newblock A complexity analysis of the F4 Gr\"obner basis algorithm with tracer data.
\newblock arXiv:2603.16378, 2026.

\bibitem{KaltofenShoup}
E.~Kaltofen and V.~Shoup.
\newblock Subquadratic-time factoring of polynomials over finite fields.
\newblock \emph{Mathematics of Computation}, 67(223):1179--1197, 1998.

\bibitem{Wiedemann}
D.~Wiedemann.
\newblock Solving sparse linear equations over finite fields.
\newblock \emph{IEEE Transactions on Information Theory}, 32(1):54--62, 1986.

\bibitem{Heintz}
J.~Heintz.
\newblock Definability and fast quantifier elimination in algebraically closed fields.
\newblock \emph{Theoretical Computer Science}, 24:239--277, 1983.

\bibitem{LachaudRolland}
G.~Lachaud and R.~Rolland.
\newblock On the number of points of algebraic sets over finite fields.
\newblock \emph{Journal of Pure and Applied Algebra}, 219(11):5117--5136, 2015.


\bibitem{Rabin1980}
M.~O.~Rabin.
\newblock Probabilistic algorithms in finite fields.
\newblock \emph{SIAM Journal on Computing}, 9(2):273--280, 1980.
\newblock doi:10.1137/0209024.

\bibitem{SafeySchost2018}
M.~Safey El Din and E.~Schost.
\newblock Bit complexity for multi-homogeneous polynomial system solving: Application to polynomial minimization.
\newblock arXiv:1605.07433v2, 2017; manuscript revised 6 November 2018.

\bibitem{Schost2003}
E.~Schost.
\newblock Computing parametric geometric resolutions.
\newblock \emph{Applicable Algebra in Engineering, Communication and Computing}, 13(5), 2002.
\newblock doi:10.1007/s00200-002-0109-x.

\bibitem{QRUOVResponse2026}
S.~Orihara, on behalf of the QR-UOV team.
\newblock Round 2 (Additional Signatures) official comment: QR-UOV.
\newblock NIST PQC Forum, 20 January 2026.

\bibitem{Hashimoto2023}
Y.~Hashimoto.
\newblock An improvement of algorithms to solve under-defined systems of multivariate quadratic equations.
\newblock \emph{JSIAM Letters}, 15:53--56, 2023.
\newblock doi:10.14495/jsiaml.15.53.

\bibitem{ThomaeWolf2012}
E.~Thomae and C.~A.~Wolf.
\newblock Solving underdetermined systems of multivariate quadratic equations revisited.
\newblock In \emph{PKC 2012}, LNCS~7293, pp.~156--171, 2012.
\newblock doi:10.1007/978-3-642-30057-8\_10.

\bibitem{FurueUnderdetermined2021}
H.~Furue, S.~Nakamura, and T.~Takagi.
\newblock Improving Thomae--Wolf algorithm for solving underdetermined multivariate quadratic polynomial problem.
\newblock In \emph{PQCrypto 2021}, LNCS~12841, pp.~65--78, 2021.
\newblock doi:10.1007/978-3-030-81293-5\_4.

\bibitem{Chevalley1935}
C.~Chevalley.
\newblock D\'emonstration d'une hypoth\`ese de M.~Artin.
\newblock \emph{Abhandlungen aus dem Mathematischen Seminar der Universit\"at Hamburg}, 11:73--75, 1935.

\bibitem{Warning1935}
E.~Warning.
\newblock Bemerkung zur vorstehenden Arbeit von Herrn Chevalley.
\newblock \emph{Abhandlungen aus dem Mathematischen Seminar der Universit\"at Hamburg}, 11:76--83, 1935.

\bibitem{JinEtAl2026}
Y.~Jin, Y.~Pan, X.~He, B.~Gong, and J.~Ding.
\newblock Security analysis on UOV families with odd characteristics: Using symmetric algebra.
\newblock Cryptology ePrint Archive, Paper 2025/1137; in \emph{PKC 2026, Part I}, LNCS~16551, pp.~428--454, 2026.

\end{thebibliography}

\end{document}
